\documentclass[aspectratio=169,10pt]{beamer}
\usetheme{Madrid}
\usecolortheme{default}
\usepackage[utf8]{inputenc}
\usepackage[T1]{fontenc}
\usepackage{lmodern}
\usepackage{amsmath,amssymb}
\usepackage{enumitem}
\usepackage{xcolor}
\setbeamertemplate{navigation symbols}{}
\setbeamertemplate{footline}[frame number]
\setbeamertemplate{itemize items}[circle]
\setbeamerfont{frametitle}{size=\large,series=\bfseries}
\title{CIA Part 2 -- Next 100 Hard Questions}
\subtitle{Engagement Planning, Information Gathering/Analysis/Evaluation, Supervision and Communication (with IT, cyber and AI focus)}
\author{Study Deck}
\date{\today}
\begin{document}
\begin{frame}\titlepage\end{frame}
\begin{frame}{Format and blueprint}
\small
\begin{itemize}
\item 100 questions, one per slide; the explanation follows each question (why correct, why the others are wrong).
\item Aligned to the three-section Part 2 structure (Internal Audit Engagement, May 2025 syllabus): Engagement Planning (50), Information Gathering, Analysis and Evaluation (40), Engagement Supervision and Communication (10). Confirm weights on the IIA syllabus page before you sit the exam.
\item Terminology follows the Global Internal Audit Standards (Domain V: Principles 13--15, plus oversight of engagement performance) and the Topical Requirements.
\item Many items are technical: ITGC, SOC reports, cloud shared responsibility, privacy, resilience, data analytics, continuous monitoring, AI.
\item Topics are seeded from your part5all--part7all files (engagement planning, evidence, workpapers, testing, supervision, communication), rewritten as harder scenario items using current Standards wording. Final-report, residual-risk and follow-up items are left out because they sit mainly in Part 3.
\end{itemize}
\end{frame}

%%%%%%%%%% Engagement Planning %%%%%%%%%%
% ===== Q1 =====
\begin{frame}[t]{Q1 \hfill {\scriptsize Engagement Planning}}
	\small
	\textbf{Question:} A business unit head asks internal audit to ``review the new vendor portal'' and wants fieldwork to begin Monday. No objectives or scope have been agreed. What should the auditor do FIRST?
	\begin{enumerate}[label=\Alph*.,itemsep=1pt,topsep=3pt]
			\item Reuse the scope of last year's portal review so the work can start on time.
			\item Start walkthroughs on Monday and let the objectives emerge from what is observed.
			\item Meet the engagement client and stakeholders to turn the request into documented objectives, scope, and criteria before committing to procedures.
			\item Draft the work program from a standard IT checklist and ask the client to approve it.
	\end{enumerate}
\end{frame}
\begin{frame}[t]{Q1 -- Explanation}
	\footnotesize
	\textbf{Correct: C.} Objectives, scope, and criteria drive every later planning step, including the risk assessment and work program. \\[5pt]
\textbf{A wrong:} Last year's scope may not reflect current risks or changes to the portal. \\[3pt]
\textbf{B wrong:} Procedures without agreed objectives invite scope drift and unfocused evidence. \\[3pt]
\textbf{D wrong:} Procedures chosen before objectives are set cannot be tied to what the engagement must conclude on. \\[3pt]
	\vfill {\tiny Part 2, Set 1}
\end{frame}
% ===== Q2 =====
\begin{frame}[t]{Q2 \hfill {\scriptsize Engagement Planning}}
	\small
	\textbf{Question:} During fieldwork on accounts payable, the team finds that a related supplier-onboarding process carries higher risk than anything in the approved scope. What is the BEST response?
	\begin{enumerate}[label=\Alph*.,itemsep=1pt,topsep=3pt]
			\item Assess whether expanding scope is warranted, obtain the required approval, and update the work program, resources, and documentation.
			\item Quietly extend testing to onboarding to avoid delaying the engagement.
			\item Cancel the engagement and restart planning from scratch.
			\item Leave onboarding out because it is outside the approved scope.
	\end{enumerate}
\end{frame}
\begin{frame}[t]{Q2 -- Explanation}
	\footnotesize
	\textbf{Correct: A.} Scope changes are legitimate when risk warrants them, but they must be approved and reflected in the plan. \\[5pt]
\textbf{B wrong:} Unapproved scope changes undermine supervision, agreed communication, and resource planning. \\[3pt]
\textbf{C wrong:} Restarting is disproportionate; a controlled scope change is sufficient. \\[3pt]
\textbf{D wrong:} Ignoring a known higher risk leaves stakeholders with a false sense of coverage. \\[3pt]
	\vfill {\tiny Part 2, Set 1}
\end{frame}
% ===== Q3 =====
\begin{frame}[t]{Q3 \hfill {\scriptsize Engagement Planning}}
	\small
	\textbf{Question:} An auditor is planning an engagement on a new customer-data sharing process for which the organization has no written policy. Which approach to evaluation criteria is BEST?
	\begin{enumerate}[label=\Alph*.,itemsep=1pt,topsep=3pt]
			\item Treat the process owner's current practice as the criterion.
			\item Defer the choice of criteria until findings are being drafted.
			\item Use the auditor's own view of good practice without discussing it with anyone.
			\item Identify suitable criteria from applicable laws, regulations, and recognized frameworks, and agree them with management where possible.
	\end{enumerate}
\end{frame}
\begin{frame}[t]{Q3 -- Explanation}
	\footnotesize
	\textbf{Correct: D.} Criteria must be relevant, authoritative enough to defend findings, and identified before evaluation begins. \\[5pt]
\textbf{A wrong:} Using the current practice as the standard is circular; no gap could ever be found. \\[3pt]
\textbf{B wrong:} Without criteria, there is nothing to measure the condition against during fieldwork. \\[3pt]
\textbf{C wrong:} Unilateral personal criteria are hard to defend and invite disputes over findings. \\[3pt]
	\vfill {\tiny Part 2, Set 1}
\end{frame}
% ===== Q4 =====
\begin{frame}[t]{Q4 \hfill {\scriptsize Engagement Planning}}
	\small
	\textbf{Question:} Management asks that a high-risk contract-management area be excluded from an engagement that has already been scoped to include it. The auditor believes the exclusion is not justified. What is the MOST appropriate action?
	\begin{enumerate}[label=\Alph*.,itemsep=1pt,topsep=3pt]
			\item Withdraw from the engagement immediately without telling anyone.
			\item Accept the exclusion and issue a conclusion covering the whole process.
			\item Escalate to the supervisor and CAE so the limitation is discussed with management and, if unresolved, communicated to the board; document the effect on conclusions.
			\item Cover the area through management inquiry alone and conclude as originally planned.
	\end{enumerate}
\end{frame}
\begin{frame}[t]{Q4 -- Explanation}
	\footnotesize
	\textbf{Correct: C.} A scope limitation that affects the ability to conclude must be escalated and its effect on the conclusion made visible. \\[5pt]
\textbf{A wrong:} Unilateral withdrawal bypasses the escalation path and leaves the limitation unresolved. \\[3pt]
\textbf{B wrong:} A conclusion that covers an excluded area overstates assurance. \\[3pt]
\textbf{D wrong:} Inquiry alone is not an adequate substitute for the planned procedures. \\[3pt]
	\vfill {\tiny Part 2, Set 1}
\end{frame}
% ===== Q5 =====
\begin{frame}[t]{Q5 \hfill {\scriptsize Engagement Planning}}
	\small
	\textbf{Question:} Which statement is an engagement EVALUATION CRITERION, as opposed to an objective, a scope item, or a procedure?
	\begin{enumerate}[label=\Alph*.,itemsep=1pt,topsep=3pt]
			\item Select 60 payroll master-data changes and inspect approval evidence.
			\item Payroll processing at the Ohio and Texas centers, July through December.
			\item Payroll master-data changes require independent approval under the HR control standard.
			\item Assess whether payroll master-data changes are appropriately approved.
	\end{enumerate}
\end{frame}
\begin{frame}[t]{Q5 -- Explanation}
	\footnotesize
	\textbf{Correct: C.} A criterion is the benchmark that defines what ``should be'' and against which the condition is evaluated. \\[5pt]
\textbf{A wrong:} This is a procedure in the work program. \\[3pt]
\textbf{B wrong:} This is scope: it sets boundaries in locations and period. \\[3pt]
\textbf{D wrong:} This is an objective: it states what the engagement intends to determine. \\[3pt]
	\vfill {\tiny Part 2, Set 1}
\end{frame}
% ===== Q6 =====
\begin{frame}[t]{Q6 \hfill {\scriptsize Engagement Planning}}
	\small
	\textbf{Question:} A company's 14 regional finance teams run differently configured instances of the same ERP. How should this structure affect the engagement risk assessment?
	\begin{enumerate}[label=\Alph*.,itemsep=1pt,topsep=3pt]
			\item Test headquarters only and extend the conclusion to every region.
			\item Assume central policies operate uniformly across regions.
			\item Reduce testing because regional teams will self-correct.
			\item Assess risk at each significant location, considering differences in control design, and avoid extending one region's results to all.
	\end{enumerate}
\end{frame}
\begin{frame}[t]{Q6 -- Explanation}
	\footnotesize
	\textbf{Correct: D.} Decentralized structures increase the risk of inconsistent control operation, so risk is assessed where the work is done. \\[5pt]
\textbf{A wrong:} Extrapolating from one location ignores configuration and operating differences. \\[3pt]
\textbf{B wrong:} Uniform operation is an assumption that must be tested, not presumed. \\[3pt]
\textbf{C wrong:} Self-correction is not evidence; it is a hope that needs to be verified. \\[3pt]
	\vfill {\tiny Part 2, Set 1}
\end{frame}
% ===== Q7 =====
\begin{frame}[t]{Q7 \hfill {\scriptsize Engagement Planning}}
	\small
	\textbf{Question:} In pre-engagement interviews, several managers say targets are met ``whatever it takes,'' and a few mention being asked to adjust figures. How should planning respond?
	\begin{enumerate}[label=\Alph*.,itemsep=1pt,topsep=3pt]
			\item Treat culture as outside the engagement-level risk assessment.
			\item Raise the assessed risk of manipulation and management override, and design less predictable procedures with less reliance on inquiry and self-assessment.
			\item Reduce sample sizes because pressure tends to improve control discipline.
			\item Rely more on management representations because leadership clearly cares about results.
	\end{enumerate}
\end{frame}
\begin{frame}[t]{Q7 -- Explanation}
	\footnotesize
	\textbf{Correct: B.} Tone and behaviors affect the likelihood that controls are bypassed, so they belong in the engagement risk assessment. \\[5pt]
\textbf{A wrong:} Culture directly influences fraud and override risk and cannot be excluded. \\[3pt]
\textbf{C wrong:} Pressure increases, rather than reduces, the incentive to bypass controls. \\[3pt]
\textbf{D wrong:} Representations from a pressured management team are the weakest basis for reliance. \\[3pt]
	\vfill {\tiny Part 2, Set 1}
\end{frame}
% ===== Q8 =====
\begin{frame}[t]{Q8 \hfill {\scriptsize Engagement Planning}}
	\small
	\textbf{Question:} The team plans to audit a digital transformation using two-week sprints. Which planning decision BEST keeps the agile approach within the Standards?
	\begin{enumerate}[label=\Alph*.,itemsep=1pt,topsep=3pt]
			\item Replace the engagement risk assessment with the sprint backlog.
			\item Waive supervisory review for sprints shorter than a month.
			\item Complete the documentation after the final sprint to preserve momentum.
			\item Define objectives, minimum documentation, and supervisory review points for each sprint so evidence and oversight requirements are still met.
	\end{enumerate}
\end{frame}
\begin{frame}[t]{Q8 -- Explanation}
	\footnotesize
	\textbf{Correct: D.} Agile changes how work is delivered, not the requirements for risk-based planning, evidence, and supervision. \\[5pt]
\textbf{A wrong:} A backlog prioritizes tasks; it does not assess risk. \\[3pt]
\textbf{B wrong:} Supervision is not scaled away by sprint length. \\[3pt]
\textbf{C wrong:} Deferred documentation weakens evidence and review during the work itself. \\[3pt]
	\vfill {\tiny Part 2, Set 1}
\end{frame}
% ===== Q9 =====
\begin{frame}[t]{Q9 \hfill {\scriptsize Engagement Planning}}
	\small
	\textbf{Question:} Auditors will review a shared-services center on another continent entirely remotely. Which planning consideration is MOST important?
	\begin{enumerate}[label=\Alph*.,itemsep=1pt,topsep=3pt]
			\item Drop walkthroughs because the work is remote.
			\item Plan alternatives to direct observation, such as live walkthroughs and data pulled directly from systems, and address data security and time-zone constraints.
			\item Treat scanned documents e-mailed by the process owner as equally reliable as direct system extracts.
			\item Keep the on-site work program unchanged.
	\end{enumerate}
\end{frame}
\begin{frame}[t]{Q9 -- Explanation}
	\footnotesize
	\textbf{Correct: B.} Remote work changes the reliability of evidence and the means of obtaining it, so the program must be adapted. \\[5pt]
\textbf{A wrong:} Walkthroughs remain key to understanding design and can be done by video and screen share. \\[3pt]
\textbf{C wrong:} Documents selected and sent by the auditee are less reliable than evidence obtained directly. \\[3pt]
\textbf{D wrong:} An unchanged program ignores the effect of remote conditions on evidence. \\[3pt]
	\vfill {\tiny Part 2, Set 1}
\end{frame}
% ===== Q10 =====
\begin{frame}[t]{Q10 \hfill {\scriptsize Engagement Planning}}
	\small
	\textbf{Question:} A financial, operational, and IT team is combined in one integrated engagement. What is the key planning requirement?
	\begin{enumerate}[label=\Alph*.,itemsep=1pt,topsep=3pt]
			\item Run separate work programs and merge conclusions at the end without reconciling them.
			\item Coordinate objectives, scope, and resources across disciplines so process and supporting-system risks are assessed together.
			\item Limit scope to risks that all disciplines have in common.
			\item Assign all work to the most senior discipline.
	\end{enumerate}
\end{frame}
\begin{frame}[t]{Q10 -- Explanation}
	\footnotesize
	\textbf{Correct: B.} The value of an integrated approach is a single, coherent view of process and technology risk, which requires joint planning. \\[5pt]
\textbf{A wrong:} Unreconciled programs create gaps, overlaps, and conflicting conclusions. \\[3pt]
\textbf{C wrong:} Restricting scope to common risks omits the specialized risks that integration should surface. \\[3pt]
\textbf{D wrong:} Seniority does not ensure coverage of each discipline's risks. \\[3pt]
	\vfill {\tiny Part 2, Set 1}
\end{frame}
% ===== Q11 =====
\begin{frame}[t]{Q11 \hfill {\scriptsize Engagement Planning}}
	\small
	\textbf{Question:} An engagement requires analysis of 5 million transactions, but only one auditor has analytics skills and is on leave for part of fieldwork. What is the BEST approach?
	\begin{enumerate}[label=\Alph*.,itemsep=1pt,topsep=3pt]
			\item Proceed as planned and mention the limitation only in the final report.
			\item Assign the analytics to a junior auditor without specific review.
			\item Evaluate options such as training, a specialist or co-source, or a revised approach, and discuss the resource limitation and its implications with the supervisor.
			\item Proceed with a 25-item manual sample and conclude on the full population.
	\end{enumerate}
\end{frame}
\begin{frame}[t]{Q11 -- Explanation}
	\footnotesize
	\textbf{Correct: C.} Resources must match the engagement's needs, and limitations are surfaced and managed during planning. \\[5pt]
\textbf{A wrong:} Waiting until the end deprives stakeholders of the chance to adjust scope or resources. \\[3pt]
\textbf{B wrong:} Unreviewed work by an inexperienced auditor on a key procedure is a quality risk. \\[3pt]
\textbf{D wrong:} A small manual sample cannot support a conclusion on 5 million items. \\[3pt]
	\vfill {\tiny Part 2, Set 1}
\end{frame}
% ===== Q12 =====
\begin{frame}[t]{Q12 \hfill {\scriptsize Engagement Planning}}
	\small
	\textbf{Question:} An engagement on a cloud migration is assigned to a team with no cloud experience. Which action BEST addresses proficiency?
	\begin{enumerate}[label=\Alph*.,itemsep=1pt,topsep=3pt]
			\item Perform the engagement and note the skills gap in the report.
			\item Arrange training, co-sourcing, or specialist support before relevant work starts, with supervision matched to the team's skills.
			\item Ask the cloud vendor's staff to perform the technical tests.
			\item Proceed on management's explanations of how the cloud is configured.
	\end{enumerate}
\end{frame}
\begin{frame}[t]{Q12 -- Explanation}
	\footnotesize
	\textbf{Correct: B.} Resources must collectively have the competencies needed; gaps are closed before the work, not disclosed after. \\[5pt]
\textbf{A wrong:} Disclosing a gap afterwards does not cure unreliable work. \\[3pt]
\textbf{C wrong:} Using the vendor's staff to test the vendor's environment impairs independence and reliability. \\[3pt]
\textbf{D wrong:} Explanations are not evidence and do not build auditor competence. \\[3pt]
	\vfill {\tiny Part 2, Set 1}
\end{frame}
% ===== Q13 =====
\begin{frame}[t]{Q13 \hfill {\scriptsize Engagement Planning}}
	\small
	\textbf{Question:} Which characteristic BEST shows that an engagement work program is adequate?
	\begin{enumerate}[label=\Alph*.,itemsep=1pt,topsep=3pt]
			\item Procedures are linked to each objective and risk, specify expected evidence and sampling approach, and are approved before fieldwork and updated with documented reasons.
			\item It repeats last year's steps for consistency.
			\item It is deliberately broad so auditors can act on judgment without documenting changes.
			\item It is finalized after fieldwork so it matches what was done.
	\end{enumerate}
\end{frame}
\begin{frame}[t]{Q13 -- Explanation}
	\footnotesize
	\textbf{Correct: A.} A work program ties procedures to objectives and risks and is controlled through approval and documented change. \\[5pt]
\textbf{B wrong:} Repeating prior steps does not respond to current risk or change. \\[3pt]
\textbf{C wrong:} Undocumented discretion prevents supervision and reperformance. \\[3pt]
\textbf{D wrong:} A program written after the work is a record of activity, not a plan. \\[3pt]
	\vfill {\tiny Part 2, Set 1}
\end{frame}
% ===== Q14 =====
\begin{frame}[t]{Q14 \hfill {\scriptsize Engagement Planning}}
	\small
	\textbf{Question:} Before fieldwork begins, which information should the auditor communicate to the engagement client?
	\begin{enumerate}[label=\Alph*.,itemsep=1pt,topsep=3pt]
			\item Preliminary conclusions formed from the risk assessment.
			\item The planned performance appraisal criteria for audit staff.
			\item The exact items that will be sampled in each test.
			\item The engagement's objectives, scope, timing, and what the client is expected to provide, such as access to people and data.
	\end{enumerate}
\end{frame}
\begin{frame}[t]{Q14 -- Explanation}
	\footnotesize
	\textbf{Correct: D.} Planning communication sets shared expectations about objectives, scope, timing, and cooperation. \\[5pt]
\textbf{A wrong:} Conclusions cannot be formed before evidence is gathered. \\[3pt]
\textbf{B wrong:} Internal staff evaluation criteria are not part of engagement communication. \\[3pt]
\textbf{C wrong:} Advance notice of samples allows evidence to be altered or prepared. \\[3pt]
	\vfill {\tiny Part 2, Set 1}
\end{frame}
% ===== Q15 =====
\begin{frame}[t]{Q15 \hfill {\scriptsize Engagement Planning}}
	\small
	\textbf{Question:} An auditor managed the vendor-onboarding process until eight months ago and is now assigned to audit it. What is the MOST appropriate action?
	\begin{enumerate}[label=\Alph*.,itemsep=1pt,topsep=3pt]
			\item Disclose the prior responsibility to the supervisor; the auditor should not provide assurance on that process, so assignments are changed.
			\item Proceed but document a statement that objectivity was maintained.
			\item Audit only the sections the auditor did not personally perform.
			\item Proceed, since cooling-off rules apply only to the CAE.
	\end{enumerate}
\end{frame}
\begin{frame}[t]{Q15 -- Explanation}
	\footnotesize
	\textbf{Correct: A.} Objectivity is presumed to be impaired when assurance covers an activity the auditor was responsible for within the prior 12 months. \\[5pt]
\textbf{B wrong:} A self-declaration does not remove a presumed impairment. \\[3pt]
\textbf{C wrong:} Partial exclusion does not address the self-review threat to the process as a whole. \\[3pt]
\textbf{D wrong:} The safeguard applies to internal auditors generally, not just the CAE. \\[3pt]
	\vfill {\tiny Part 2, Set 1}
\end{frame}
% ===== Q16 =====
\begin{frame}[t]{Q16 \hfill {\scriptsize Engagement Planning}}
	\small
	\textbf{Question:} Planning to rely on a cloud payroll provider's SOC 2 Type II report, what should the auditor do FIRST?
	\begin{enumerate}[label=\Alph*.,itemsep=1pt,topsep=3pt]
			\item Treat a Type I report as evidence that controls operated effectively over the year.
			\item Rely on it once the provider's management confirms it is accurate.
			\item Evaluate the report's scope, period, service auditor competence and objectivity, relevant controls, exceptions, and the user-entity controls the organization must operate.
			\item Rely on it automatically because the opinion is unqualified.
	\end{enumerate}
\end{frame}
\begin{frame}[t]{Q16 -- Explanation}
	\footnotesize
	\textbf{Correct: C.} Reliance on another party's work requires judging whether it covers the relevant controls and period and what remains the user's responsibility. \\[5pt]
\textbf{A wrong:} A Type I report addresses design at a point in time, not operating effectiveness. \\[3pt]
\textbf{B wrong:} Management's confirmation is not an independent evaluation of the report. \\[3pt]
\textbf{D wrong:} An unqualified opinion on irrelevant controls or a different period adds little. \\[3pt]
	\vfill {\tiny Part 2, Set 1}
\end{frame}
% ===== Q17 =====
\begin{frame}[t]{Q17 \hfill {\scriptsize Engagement Planning}}
	\small
	\textbf{Question:} An engagement will assess cybersecurity governance. How should the applicable Topical Requirement be treated?
	\begin{enumerate}[label=\Alph*.,itemsep=1pt,topsep=3pt]
			\item Drop the requirement if the process owner calls the area operational.
			\item Treat it as optional supplemental guidance.
			\item Apply it only to engagements labeled as IT audits.
			\item Apply it as a minimum baseline for governance, risk management, and control, excluding an individual element only with documented rationale.
	\end{enumerate}
\end{frame}
\begin{frame}[t]{Q17 -- Explanation}
	\footnotesize
	\textbf{Correct: D.} Topical Requirements are mandatory where they apply to the subject of an assurance engagement. \\[5pt]
\textbf{A wrong:} Applicability turns on the subject matter, not on how the owner labels the area. \\[3pt]
\textbf{B wrong:} They are part of the mandatory elements of the framework, not optional guidance. \\[3pt]
\textbf{C wrong:} Applicability is not limited to engagements with an IT label. \\[3pt]
	\vfill {\tiny Part 2, Set 1}
\end{frame}
% ===== Q18 =====
\begin{frame}[t]{Q18 \hfill {\scriptsize Engagement Planning}}
	\small
	\textbf{Question:} Before relying on an automated three-way match in an accounts payable engagement, what should the auditor do?
	\begin{enumerate}[label=\Alph*.,itemsep=1pt,topsep=3pt]
			\item Review the IT budget for the application.
			\item Obtain assurance that general IT controls over the application, such as access, change management, and operations, are effective, then test the match configuration.
			\item Interview the developer about how the match was designed.
			\item Run a few invoices through the system and conclude if they match.
	\end{enumerate}
\end{frame}
\begin{frame}[t]{Q18 -- Explanation}
	\footnotesize
	\textbf{Correct: B.} Automated controls are only as reliable as the IT environment that protects them from unauthorized change. \\[5pt]
\textbf{A wrong:} Budgets say nothing about control reliability. \\[3pt]
\textbf{C wrong:} Interviews give design information but not evidence of effective operation. \\[3pt]
\textbf{D wrong:} A few passing items do not show the control is protected and consistently applied. \\[3pt]
	\vfill {\tiny Part 2, Set 1}
\end{frame}
% ===== Q19 =====
\begin{frame}[t]{Q19 \hfill {\scriptsize Engagement Planning}}
	\small
	\textbf{Question:} Which objective BEST addresses resilience in an engagement on disaster recovery for payment processing?
	\begin{enumerate}[label=\Alph*.,itemsep=1pt,topsep=3pt]
			\item Determine whether a disaster recovery policy exists and has been approved.
			\item Determine whether the data center has redundant power.
			\item Determine whether tested recovery capabilities meet the recovery time and recovery point objectives derived from the business impact analysis.
			\item Determine whether backups run on a daily schedule.
	\end{enumerate}
\end{frame}
\begin{frame}[t]{Q19 -- Explanation}
	\footnotesize
	\textbf{Correct: C.} The meaningful question is whether the organization can actually recover within what the business requires. \\[5pt]
\textbf{A wrong:} A policy shows intent, not capability. \\[3pt]
\textbf{B wrong:} Redundant power is one component and does not demonstrate recovery. \\[3pt]
\textbf{D wrong:} Backup schedules do not show data can be restored within required time. \\[3pt]
	\vfill {\tiny Part 2, Set 1}
\end{frame}
% ===== Q20 =====
\begin{frame}[t]{Q20 \hfill {\scriptsize Engagement Planning}}
	\small
	\textbf{Question:} An engagement will review customer data sharing with third-party processors under privacy law. Which planning element is MOST important?
	\begin{enumerate}[label=\Alph*.,itemsep=1pt,topsep=3pt]
			\item Defer data mapping until fieldwork is complete.
			\item Rely on the legal department's assertion of compliance.
			\item Identify the applicable legal obligations as criteria and map the data flows and processors to define scope.
			\item Limit scope to IT security controls on internally owned systems.
	\end{enumerate}
\end{frame}
\begin{frame}[t]{Q20 -- Explanation}
	\footnotesize
	\textbf{Correct: C.} Criteria come from legal duties, and scope must follow where personal data actually goes, including to third parties. \\[5pt]
\textbf{A wrong:} Without mapping, scope cannot be set responsibly. \\[3pt]
\textbf{B wrong:} An assertion is not evidence and does not define criteria or scope. \\[3pt]
\textbf{D wrong:} Privacy compliance extends beyond security, and processors hold the data off-system. \\[3pt]
	\vfill {\tiny Part 2, Set 1}
\end{frame}
% ===== Q21 =====
\begin{frame}[t]{Q21 \hfill {\scriptsize Engagement Planning}}
	\small
	\textbf{Question:} Planning analytics show days sales outstanding up 30\% while revenue is up 12\% and returns have fallen. What is the MOST reasonable planning response?
	\begin{enumerate}[label=\Alph*.,itemsep=1pt,topsep=3pt]
			\item Conclude there is no risk because revenue is growing.
			\item Treat it as an indicator of possible revenue-recognition, cutoff, or collectability issues and add targeted procedures.
			\item Conclude collections have improved.
			\item Treat it as an accounts payable issue.
	\end{enumerate}
\end{frame}
\begin{frame}[t]{Q21 -- Explanation}
	\footnotesize
	\textbf{Correct: B.} Receivables growing faster than sales is a classic indicator that revenue timing or valuation may be misstated. \\[5pt]
\textbf{A wrong:} Growing revenue can be exactly the setting for aggressive recognition. \\[3pt]
\textbf{C wrong:} Rising days outstanding indicates slower, not faster, collection. \\[3pt]
\textbf{D wrong:} Receivables relate to customers, not payables. \\[3pt]
	\vfill {\tiny Part 2, Set 1}
\end{frame}
% ===== Q22 =====
\begin{frame}[t]{Q22 \hfill {\scriptsize Engagement Planning}}
	\small
	\textbf{Question:} In planning a vendor master-file engagement, which objective is MOST appropriate?
	\begin{enumerate}[label=\Alph*.,itemsep=1pt,topsep=3pt]
			\item Evaluate the storage cost of the vendor database.
			\item Evaluate whether controls prevent and detect unauthorized creation or change of vendors that could enable fictitious or duplicate payments.
			\item Evaluate whether vendor names follow a consistent naming format.
			\item Evaluate how many vendors are held in the master file.
	\end{enumerate}
\end{frame}
\begin{frame}[t]{Q22 -- Explanation}
	\footnotesize
	\textbf{Correct: B.} Objectives should address the risks that matter, here payment fraud and error through master data. \\[5pt]
\textbf{A wrong:} Storage cost is an efficiency matter, not a control risk. \\[3pt]
\textbf{C wrong:} Formatting matters for hygiene, but it is not the key risk. \\[3pt]
\textbf{D wrong:} Volume alone does not indicate control failure. \\[3pt]
	\vfill {\tiny Part 2, Set 1}
\end{frame}
% ===== Q23 =====
\begin{frame}[t]{Q23 \hfill {\scriptsize Engagement Planning}}
	\small
	\textbf{Question:} Which approach to prioritizing risks and controls in an engagement risk assessment is BEST?
	\begin{enumerate}[label=\Alph*.,itemsep=1pt,topsep=3pt]
			\item Rank areas by the number of findings in last year's report only.
			\item Rank areas by how familiar the auditors are with them.
			\item Rank controls by how cheap they are to test.
			\item Rate inherent risk by impact and likelihood, identify key controls that mitigate it, and decide reliance based on design and testing.
	\end{enumerate}
\end{frame}
\begin{frame}[t]{Q23 -- Explanation}
	\footnotesize
	\textbf{Correct: D.} A risk-based approach links risk, mitigating controls, and the extent of testing. \\[5pt]
\textbf{A wrong:} Prior findings are one input, not a complete picture of current risk. \\[3pt]
\textbf{B wrong:} Familiarity is a convenience, not a risk measure. \\[3pt]
\textbf{C wrong:} Cost of testing may affect method but should not set priority of risk. \\[3pt]
	\vfill {\tiny Part 2, Set 1}
\end{frame}
% ===== Q24 =====
\begin{frame}[t]{Q24 \hfill {\scriptsize Engagement Planning}}
	\small
	\textbf{Question:} A company has just launched a generative AI chatbot to handle customer claims. What is the BEST way to approach engagement risk assessment?
	\begin{enumerate}[label=\Alph*.,itemsep=1pt,topsep=3pt]
			\item Identify AI-specific risks, such as accuracy, bias, privacy, third-party model dependence, and human oversight, and link them to objectives and criteria such as the organization's AI governance policy.
			\item Rely on the vendor's published model documentation.
			\item Treat it as an ordinary application and limit scope to general IT controls.
			\item Wait until regulators finalize AI rules before planning.
	\end{enumerate}
\end{frame}
\begin{frame}[t]{Q24 -- Explanation}
	\footnotesize
	\textbf{Correct: A.} New technology changes the risk profile, so planning must reflect the specific risks it introduces. \\[5pt]
\textbf{B wrong:} Vendor documentation is the vendor's assertion and not independent assurance. \\[3pt]
\textbf{C wrong:} General IT controls do not cover model behavior or outcomes. \\[3pt]
\textbf{D wrong:} Governance, risk, and control exist now, and delay leaves exposure unexamined. \\[3pt]
	\vfill {\tiny Part 2, Set 1}
\end{frame}
% ===== Q25 =====
\begin{frame}[t]{Q25 \hfill {\scriptsize Engagement Planning}}
	\small
	\textbf{Question:} A hotline tip says managers at one branch approve their own expense claims. How should this affect planning of an expense engagement?
	\begin{enumerate}[label=\Alph*.,itemsep=1pt,topsep=3pt]
			\item Raise the assessed fraud and override risk and design procedures such as analytics on self-approvals and unannounced testing.
			\item Ask the branch manager to explain before deciding on any testing.
			\item Refer it to another team and leave the branch out of the engagement.
			\item Limit testing to low-value claims.
	\end{enumerate}
\end{frame}
\begin{frame}[t]{Q25 -- Explanation}
	\footnotesize
	\textbf{Correct: A.} Planning must consider the probability of fraud and noncompliance and shape procedures accordingly. \\[5pt]
\textbf{B wrong:} Tipping off the implicated manager risks evidence being altered. \\[3pt]
\textbf{C wrong:} Referral may be proper for investigation but does not remove the planning obligation for this engagement. \\[3pt]
\textbf{D wrong:} Low-value testing ignores where abuse is likely to concentrate. \\[3pt]
	\vfill {\tiny Part 2, Set 1}
\end{frame}
% ===== Q26 =====
\begin{frame}[t]{Q26 \hfill {\scriptsize Engagement Planning}}
	\small
	\textbf{Question:} After planning is complete, an engagement's budget is cut by 30\%. What is the BEST response?
	\begin{enumerate}[label=\Alph*.,itemsep=1pt,topsep=3pt]
			\item Keep the work program and shorten supervision instead.
			\item Keep scope and reduce every sample uniformly without telling anyone.
			\item Reassess risk-based priorities, adjust scope or procedures with approval, and document the effect on the level of assurance.
			\item Omit documentation to save time.
	\end{enumerate}
\end{frame}
\begin{frame}[t]{Q26 -- Explanation}
	\footnotesize
	\textbf{Correct: C.} Resource constraints require a conscious, approved trade-off that stakeholders can understand. \\[5pt]
\textbf{A wrong:} Cutting supervision lowers quality while the effect stays invisible. \\[3pt]
\textbf{B wrong:} Hidden reductions leave stakeholders believing assurance is unchanged. \\[3pt]
\textbf{D wrong:} Documentation is required to support conclusions. \\[3pt]
	\vfill {\tiny Part 2, Set 1}
\end{frame}
% ===== Q27 =====
\begin{frame}[t]{Q27 \hfill {\scriptsize Engagement Planning}}
	\small
	\textbf{Question:} Purchasing-card analytics show 3\% of transactions sit just below the approval threshold. How should this be used in planning?
	\begin{enumerate}[label=\Alph*.,itemsep=1pt,topsep=3pt]
			\item Ignore it and sample uniformly across all transactions.
			\item Notify cardholders so they can correct their behavior.
			\item Target the work program on threshold avoidance, including split purchases, and select items from the flagged group.
			\item Conclude the approval control is effective.
	\end{enumerate}
\end{frame}
\begin{frame}[t]{Q27 -- Explanation}
	\footnotesize
	\textbf{Correct: C.} Analytics results that point to a specific risk should shape scope and procedures. \\[5pt]
\textbf{A wrong:} Uniform sampling dilutes attention on where the risk is concentrated. \\[3pt]
\textbf{B wrong:} Pre-notification would undermine the evidence and any later investigation. \\[3pt]
\textbf{D wrong:} Clustering below a threshold suggests the control is being circumvented. \\[3pt]
	\vfill {\tiny Part 2, Set 1}
\end{frame}
% ===== Q28 =====
\begin{frame}[t]{Q28 \hfill {\scriptsize Engagement Planning}}
	\small
	\textbf{Question:} The objective is to estimate the percentage of paid invoices in error across the entire population with measurable confidence. Which sampling approach is BEST?
	\begin{enumerate}[label=\Alph*.,itemsep=1pt,topsep=3pt]
			\item A judgmental sample of the 40 largest invoices.
			\item The first 30 invoices of each month.
			\item A stratified sample of high-value invoices only.
			\item Statistical random sampling sized for the desired confidence and tolerable error.
	\end{enumerate}
\end{frame}
\begin{frame}[t]{Q28 -- Explanation}
	\footnotesize
	\textbf{Correct: D.} Quantified conclusions about a population require a sample whose results can be projected statistically. \\[5pt]
\textbf{A wrong:} Judgmental selections cannot be projected to the population. \\[3pt]
\textbf{B wrong:} A non-random selection is biased and cannot support a measured estimate. \\[3pt]
\textbf{C wrong:} Sampling only high-value items restricts conclusions to that stratum. \\[3pt]
	\vfill {\tiny Part 2, Set 1}
\end{frame}
% ===== Q29 =====
\begin{frame}[t]{Q29 \hfill {\scriptsize Engagement Planning}}
	\small
	\textbf{Question:} The auditor lowers the tolerable deviation rate for a key control test. What is the effect on sample size, other factors unchanged?
	\begin{enumerate}[label=\Alph*.,itemsep=1pt,topsep=3pt]
			\item The sample size depends only on population size.
			\item The required sample size decreases.
			\item The required sample size increases.
			\item The sample size is unchanged.
	\end{enumerate}
\end{frame}
\begin{frame}[t]{Q29 -- Explanation}
	\footnotesize
	\textbf{Correct: C.} A tighter tolerance requires more evidence to be confident that deviations are below the threshold. \\[5pt]
\textbf{A wrong:} Population size has a limited effect compared with tolerance, risk, and confidence. \\[3pt]
\textbf{B wrong:} A smaller tolerance cannot be demonstrated with less evidence. \\[3pt]
\textbf{D wrong:} Tolerance is a direct driver of sample size. \\[3pt]
	\vfill {\tiny Part 2, Set 1}
\end{frame}
% ===== Q30 =====
\begin{frame}[t]{Q30 \hfill {\scriptsize Engagement Planning}}
	\small
	\textbf{Question:} A key approval workflow was reconfigured on August 1; the audit period is January to December. Which testing approach is BEST?
	\begin{enumerate}[label=\Alph*.,itemsep=1pt,topsep=3pt]
			\item Rely on a screenshot of the current configuration.
			\item Test only items after the change.
			\item Treat pre- and post-change periods as separate populations and test the control as configured in each.
			\item Select a random sample across the year and ignore the change.
	\end{enumerate}
\end{frame}
\begin{frame}[t]{Q30 -- Explanation}
	\footnotesize
	\textbf{Correct: C.} A control that changes during the period must be evaluated as it operated throughout the period. \\[5pt]
\textbf{A wrong:} A current screenshot says nothing about how the control worked before the change. \\[3pt]
\textbf{B wrong:} Testing only the later period leaves the earlier months uncovered. \\[3pt]
\textbf{D wrong:} A single pooled sample mixes two different controls and weakens conclusions. \\[3pt]
	\vfill {\tiny Part 2, Set 1}
\end{frame}
% ===== Q31 =====
\begin{frame}[t]{Q31 \hfill {\scriptsize Engagement Planning}}
	\small
	\textbf{Question:} The engagement overlaps with the external auditor's scope this year. What is the BEST approach?
	\begin{enumerate}[label=\Alph*.,itemsep=1pt,topsep=3pt]
			\item Avoid contact so independence is preserved.
			\item Let the external auditor set the engagement objectives.
			\item Adopt the external auditor's risk assessment without evaluation.
			\item Coordinate to understand their scope and timing to limit duplication, while evaluating any work relied upon.
	\end{enumerate}
\end{frame}
\begin{frame}[t]{Q31 -- Explanation}
	\footnotesize
	\textbf{Correct: D.} Coordination reduces gaps and duplication but does not replace the internal auditor's own judgment. \\[5pt]
\textbf{A wrong:} Coordination does not impair independence; isolation creates gaps and duplication. \\[3pt]
\textbf{B wrong:} The engagement's objectives are determined by internal audit's own planning. \\[3pt]
\textbf{C wrong:} Reliance requires evaluating competence, objectivity, and scope of that work. \\[3pt]
	\vfill {\tiny Part 2, Set 1}
\end{frame}
% ===== Q32 =====
\begin{frame}[t]{Q32 \hfill {\scriptsize Engagement Planning}}
	\small
	\textbf{Question:} An outsourced call center's contract includes a right-to-audit clause. What is the BEST planning approach?
	\begin{enumerate}[label=\Alph*.,itemsep=1pt,topsep=3pt]
			\item Audit only the internal vendor-management tracker.
			\item Use the clause to agree scope with the vendor, consider subcontractors and confidentiality, and plan direct evidence rather than relying only on vendor-prepared material.
			\item Exclude the vendor because it is an independent company.
			\item Rely entirely on the vendor's self-assessment questionnaire.
	\end{enumerate}
\end{frame}
\begin{frame}[t]{Q32 -- Explanation}
	\footnotesize
	\textbf{Correct: B.} Third-party risk remains the organization's, and contractual access should be used to obtain direct evidence. \\[5pt]
\textbf{A wrong:} A tracker shows oversight activity but not how the service performs. \\[3pt]
\textbf{C wrong:} Outsourcing transfers activity, not accountability for risk. \\[3pt]
\textbf{D wrong:} Self-assessments are assertions and not evidence of operation. \\[3pt]
	\vfill {\tiny Part 2, Set 1}
\end{frame}
% ===== Q33 =====
\begin{frame}[t]{Q33 \hfill {\scriptsize Engagement Planning}}
	\small
	\textbf{Question:} An auditor will extract sensitive customer data into an analytics tool. Which planning step is MOST appropriate?
	\begin{enumerate}[label=\Alph*.,itemsep=1pt,topsep=3pt]
			\item Skip anonymization since auditors are trusted.
			\item Plan extraction of only the data needed, with secure storage, approvals, and access limited to the team.
			\item Upload the data to an external AI service to speed up analysis.
			\item Copy the full population to a personal laptop for speed.
	\end{enumerate}
\end{frame}
\begin{frame}[t]{Q33 -- Explanation}
	\footnotesize
	\textbf{Correct: B.} Audit work involves handling sensitive data, so it must follow data-minimization and security principles. \\[5pt]
\textbf{A wrong:} Trust does not replace safeguards and minimization. \\[3pt]
\textbf{C wrong:} Sending data to unapproved third parties breaches confidentiality obligations. \\[3pt]
\textbf{D wrong:} Personal devices create security and confidentiality exposure. \\[3pt]
	\vfill {\tiny Part 2, Set 1}
\end{frame}
% ===== Q34 =====
\begin{frame}[t]{Q34 \hfill {\scriptsize Engagement Planning}}
	\small
	\textbf{Question:} Internal policy allows a retention period shorter than a legal minimum. Which criterion should the auditor use?
	\begin{enumerate}[label=\Alph*.,itemsep=1pt,topsep=3pt]
			\item The internal policy, because management approved it.
			\item The legal minimum, treating noncompliance with law as a finding even if the policy was followed.
			\item Whichever is less strict, to avoid disputes.
			\item The average of the policy and the law.
	\end{enumerate}
\end{frame}
\begin{frame}[t]{Q34 -- Explanation}
	\footnotesize
	\textbf{Correct: B.} Laws and regulations set the floor; policy cannot lawfully go below them. \\[5pt]
\textbf{A wrong:} Management approval cannot override legal requirements. \\[3pt]
\textbf{C wrong:} Choosing the less strict source would hide a compliance exposure. \\[3pt]
\textbf{D wrong:} Averaging has no basis in either source. \\[3pt]
	\vfill {\tiny Part 2, Set 1}
\end{frame}
% ===== Q35 =====
\begin{frame}[t]{Q35 \hfill {\scriptsize Engagement Planning}}
	\small
	\textbf{Question:} An actuarial specialist will support an engagement on reserves. Which statement about responsibility is MOST accurate?
	\begin{enumerate}[label=\Alph*.,itemsep=1pt,topsep=3pt]
			\item The specialist's report can be adopted as the conclusion without further consideration.
			\item Responsibility for the conclusion transfers to the specialist.
			\item The auditor should evaluate the specialist's competence and objectivity, understand the assumptions and scope, and remain responsible for the engagement conclusion.
			\item Evaluation of the specialist is needed only if problems arise.
	\end{enumerate}
\end{frame}
\begin{frame}[t]{Q35 -- Explanation}
	\footnotesize
	\textbf{Correct: C.} Using a specialist does not transfer accountability for the conclusion. \\[5pt]
\textbf{A wrong:} Adopting a report unexamined skips the auditor's own judgment. \\[3pt]
\textbf{B wrong:} Accountability for the engagement conclusion stays with internal audit. \\[3pt]
\textbf{D wrong:} Evaluation is done up front, since problems may not be visible later. \\[3pt]
	\vfill {\tiny Part 2, Set 1}
\end{frame}
% ===== Q36 =====
\begin{frame}[t]{Q36 \hfill {\scriptsize Engagement Planning}}
	\small
	\textbf{Question:} The prior report for the area under review noted an overdue action plan to restrict privileged access. How should this be handled when planning?
	\begin{enumerate}[label=\Alph*.,itemsep=1pt,topsep=3pt]
			\item Consider it in the risk assessment and scope, and test the current state of privileged access if it remains relevant.
			\item Repeat the old finding without testing.
			\item Exclude it because it has already been reported.
			\item Assume it is resolved when management says so.
	\end{enumerate}
\end{frame}
\begin{frame}[t]{Q36 -- Explanation}
	\footnotesize
	\textbf{Correct: A.} Known unresolved issues inform current risk and should be reflected in objectives and procedures. \\[5pt]
\textbf{B wrong:} Repeating a finding without current evidence is not supportable. \\[3pt]
\textbf{C wrong:} A reported issue can remain a live risk. \\[3pt]
\textbf{D wrong:} Management's statement is not evidence the issue has been remediated. \\[3pt]
	\vfill {\tiny Part 2, Set 1}
\end{frame}
% ===== Q37 =====
\begin{frame}[t]{Q37 \hfill {\scriptsize Engagement Planning}}
	\small
	\textbf{Question:} In an IaaS arrangement under the shared responsibility model, which control MOST likely remains the customer's responsibility?
	\begin{enumerate}[label=\Alph*.,itemsep=1pt,topsep=3pt]
			\item Patching of the underlying hypervisor.
			\item Physical security of the data center.
			\item Power and cooling of the facility.
			\item Configuring user access and patching the operating systems of the customer's virtual machines.
	\end{enumerate}
\end{frame}
\begin{frame}[t]{Q37 -- Explanation}
	\footnotesize
	\textbf{Correct: D.} In IaaS, the customer manages its own operating systems, identities, and data configuration. \\[5pt]
\textbf{A wrong:} The provider manages the virtualization layer in IaaS. \\[3pt]
\textbf{B wrong:} Physical security is the provider's responsibility. \\[3pt]
\textbf{C wrong:} Facilities are the provider's domain. \\[3pt]
	\vfill {\tiny Part 2, Set 1}
\end{frame}
% ===== Q38 =====
\begin{frame}[t]{Q38 \hfill {\scriptsize Engagement Planning}}
	\small
	\textbf{Question:} The organization holds consigned stock at three third-party warehouses. How should this affect engagement planning on inventory?
	\begin{enumerate}[label=\Alph*.,itemsep=1pt,topsep=3pt]
			\item Rely on the perpetual inventory records alone.
			\item Accept an e-mail from the warehouse manager stating quantities.
			\item Include those sites in scope and plan to observe counts or obtain direct confirmation from the warehouse operators.
			\item Exclude them because the stock is held off-site.
	\end{enumerate}
\end{frame}
\begin{frame}[t]{Q38 -- Explanation}
	\footnotesize
	\textbf{Correct: C.} Assets held by third parties still affect the organization's records and must be covered by planned evidence. \\[5pt]
\textbf{A wrong:} Internal records alone cannot confirm physical existence. \\[3pt]
\textbf{B wrong:} E-mailed statements are weak, unsolicited evidence. \\[3pt]
\textbf{D wrong:} Off-site stock is still reported and still carries risk. \\[3pt]
	\vfill {\tiny Part 2, Set 1}
\end{frame}
% ===== Q39 =====
\begin{frame}[t]{Q39 \hfill {\scriptsize Engagement Planning}}
	\small
	\textbf{Question:} Key financial-close calculations are built in unprotected spreadsheets. Which engagement objective is MOST appropriate?
	\begin{enumerate}[label=\Alph*.,itemsep=1pt,topsep=3pt]
			\item Evaluate end-user computing controls such as access, version control, formula integrity, and reconciliation of inputs to source.
			\item Confirm the preparer's statement that totals are correct.
			\item Evaluate whether the spreadsheets use consistent fonts and layout.
			\item Evaluate whether the spreadsheets are backed up to cloud storage.
	\end{enumerate}
\end{frame}
\begin{frame}[t]{Q39 -- Explanation}
	\footnotesize
	\textbf{Correct: A.} Spreadsheet risk lies in unauthorized changes, formula errors, and unreliable inputs. \\[5pt]
\textbf{B wrong:} A preparer's assurance is not independent evidence. \\[3pt]
\textbf{C wrong:} Formatting is not a risk to the numbers. \\[3pt]
\textbf{D wrong:} Backups protect availability, not accuracy. \\[3pt]
	\vfill {\tiny Part 2, Set 1}
\end{frame}
% ===== Q40 =====
\begin{frame}[t]{Q40 \hfill {\scriptsize Engagement Planning}}
	\small
	\textbf{Question:} In a risk and control matrix used to build the work program, which linkage is ESSENTIAL?
	\begin{enumerate}[label=\Alph*.,itemsep=1pt,topsep=3pt]
			\item Each test mapped to the auditor who performed it last year.
			\item Each control mapped to the cost of testing it.
			\item Each risk mapped to the person who identified it.
			\item Each key risk mapped to the key control or controls that address it and to the test procedure for each control.
	\end{enumerate}
\end{frame}
\begin{frame}[t]{Q40 -- Explanation}
	\footnotesize
	\textbf{Correct: D.} The matrix is useful only if it connects risks, controls, and procedures. \\[5pt]
\textbf{A wrong:} Prior staffing does not show coverage of risk. \\[3pt]
\textbf{B wrong:} Cost is a resourcing consideration, not a linkage that proves coverage. \\[3pt]
\textbf{C wrong:} Who identified a risk does not show whether it is addressed. \\[3pt]
	\vfill {\tiny Part 2, Set 1}
\end{frame}
% ===== Q41 =====
\begin{frame}[t]{Q41 \hfill {\scriptsize Engagement Planning}}
	\small
	\textbf{Question:} Which of the following is an example of ENGAGEMENT risk, as opposed to risk in the area under review?
	\begin{enumerate}[label=\Alph*.,itemsep=1pt,topsep=3pt]
			\item An incomplete data extract leads the team to an incorrect conclusion.
			\item The company misses a regulatory filing deadline.
			\item A key supplier becomes insolvent.
			\item The payroll team processes a duplicate payment.
	\end{enumerate}
\end{frame}
\begin{frame}[t]{Q41 -- Explanation}
	\footnotesize
	\textbf{Correct: A.} Engagement risk is the risk that the audit work itself fails to meet its objectives or reaches a wrong conclusion. \\[5pt]
\textbf{B wrong:} This is a compliance risk for the organization. \\[3pt]
\textbf{C wrong:} This is a supply-chain risk to the organization. \\[3pt]
\textbf{D wrong:} This is a process risk in the area under review. \\[3pt]
	\vfill {\tiny Part 2, Set 1}
\end{frame}
% ===== Q42 =====
\begin{frame}[t]{Q42 \hfill {\scriptsize Engagement Planning}}
	\small
	\textbf{Question:} Three areas are in scope: Area X has high inherent risk and weak controls; Area Y has high inherent risk and strong, previously tested controls; Area Z has low inherent risk. How should resources be allocated?
	\begin{enumerate}[label=\Alph*.,itemsep=1pt,topsep=3pt]
			\item The same effort as last year for each area.
			\item Most to Y because it offers the most evidence to test.
			\item Most to X, a moderate amount to Y to confirm key controls, and minimal to Z.
			\item Equal effort to all three areas.
	\end{enumerate}
\end{frame}
\begin{frame}[t]{Q42 -- Explanation}
	\footnotesize
	\textbf{Correct: C.} Effort should follow risk after considering the strength of controls. \\[5pt]
\textbf{A wrong:} Prior-year effort does not reflect current risk. \\[3pt]
\textbf{B wrong:} Strong controls reduce the need for extensive work rather than increase it. \\[3pt]
\textbf{D wrong:} Equal effort ignores differences in risk. \\[3pt]
	\vfill {\tiny Part 2, Set 1}
\end{frame}
% ===== Q43 =====
\begin{frame}[t]{Q43 \hfill {\scriptsize Engagement Planning}}
	\small
	\textbf{Question:} Planning concludes the control environment is strong. What is the effect on the engagement approach?
	\begin{enumerate}[label=\Alph*.,itemsep=1pt,topsep=3pt]
			\item The extent of some testing may be reduced, but targeted procedures remain necessary, including for management override.
			\item Transaction testing can be eliminated.
			\item Sample sizes should be increased to match the quality of the environment.
			\item Journal-entry testing can be excluded because override risk is nil.
	\end{enumerate}
\end{frame}
\begin{frame}[t]{Q43 -- Explanation}
	\footnotesize
	\textbf{Correct: A.} A strong environment supports reliance but does not prevent specific errors or override. \\[5pt]
\textbf{B wrong:} No amount of environment strength removes the need for evidence. \\[3pt]
\textbf{C wrong:} Better environments justify less, not more, testing. \\[3pt]
\textbf{D wrong:} Override risk is never eliminated by a strong environment. \\[3pt]
	\vfill {\tiny Part 2, Set 1}
\end{frame}
% ===== Q44 =====
\begin{frame}[t]{Q44 \hfill {\scriptsize Engagement Planning}}
	\small
	\textbf{Question:} Which sequence of engagement planning activities is MOST logical?
	\begin{enumerate}[label=\Alph*.,itemsep=1pt,topsep=3pt]
			\item Select samples, then define objectives and scope.
			\item Outline the report, then define criteria and assess risks.
			\item Understand the area, assess engagement risks, set objectives and scope, identify criteria, determine resources, and build the work program.
			\item Build the work program, then assess risks, then set objectives.
	\end{enumerate}
\end{frame}
\begin{frame}[t]{Q44 -- Explanation}
	\footnotesize
	\textbf{Correct: C.} Each step provides inputs to the next, with the work program following from risks, objectives, and scope. \\[5pt]
\textbf{A wrong:} Samples depend on objectives and risks and cannot come first. \\[3pt]
\textbf{B wrong:} A report outline cannot be set before there are objectives and evidence. \\[3pt]
\textbf{D wrong:} A program built before risks and objectives has nothing to be based on. \\[3pt]
	\vfill {\tiny Part 2, Set 1}
\end{frame}
% ===== Q45 =====
\begin{frame}[t]{Q45 \hfill {\scriptsize Engagement Planning}}
	\small
	\textbf{Question:} In planning a public-procurement compliance engagement, which analytic indicator BEST targets potential noncompliance?
	\begin{enumerate}[label=\Alph*.,itemsep=1pt,topsep=3pt]
			\item Average days to pay vendors.
			\item The number of active vendors.
			\item Awards to vendors sharing addresses or bank details with employees, or repeated purchases just under the competitive-bid threshold.
			\item Total spend by category.
	\end{enumerate}
\end{frame}
\begin{frame}[t]{Q45 -- Explanation}
	\footnotesize
	\textbf{Correct: C.} Red flags such as related parties and threshold avoidance point directly to bypassed procurement rules. \\[5pt]
\textbf{A wrong:} Payment speed is an efficiency measure. \\[3pt]
\textbf{B wrong:} Vendor counts do not show noncompliance. \\[3pt]
\textbf{D wrong:} Spend totals do not indicate whether rules were followed. \\[3pt]
	\vfill {\tiny Part 2, Set 1}
\end{frame}
% ===== Q46 =====
\begin{frame}[t]{Q46 \hfill {\scriptsize Engagement Planning}}
	\small
	\textbf{Question:} During fieldwork, an auditor wants to drop three planned procedures that appear low-yield. What is the MOST appropriate action?
	\begin{enumerate}[label=\Alph*.,itemsep=1pt,topsep=3pt]
			\item Drop them and move on.
			\item Obtain supervisory approval and document the rationale and the effect on coverage of objectives.
			\item Replace them without updating the work program.
			\item Keep every procedure regardless of what has been learned.
	\end{enumerate}
\end{frame}
\begin{frame}[t]{Q46 -- Explanation}
	\footnotesize
	\textbf{Correct: B.} Work programs may change as understanding improves, but changes are controlled and documented. \\[5pt]
\textbf{A wrong:} Unapproved changes remove supervisory control over coverage. \\[3pt]
\textbf{C wrong:} An outdated program no longer represents what was done. \\[3pt]
\textbf{D wrong:} Rigid adherence ignores new information about risk. \\[3pt]
	\vfill {\tiny Part 2, Set 1}
\end{frame}
% ===== Q47 =====
\begin{frame}[t]{Q47 \hfill {\scriptsize Engagement Planning}}
	\small
	\textbf{Question:} The IT team shares a penetration test report showing no critical findings. How should it be used in planning a cybersecurity engagement?
	\begin{enumerate}[label=\Alph*.,itemsep=1pt,topsep=3pt]
			\item As one input, after evaluating its scope, date, tester independence and competence, and coverage.
			\item Disregarded because IT arranged it.
			\item As full assurance, removing cyber from the scope.
			\item Accepted because the summary says ``passed.''
	\end{enumerate}
\end{frame}
\begin{frame}[t]{Q47 -- Explanation}
	\footnotesize
	\textbf{Correct: A.} Others' work can inform planning, but reliance depends on whether it covers the relevant risk. \\[5pt]
\textbf{B wrong:} Ignoring it wastes useful information that can be evaluated. \\[3pt]
\textbf{C wrong:} A single test rarely covers governance, process, and all assets. \\[3pt]
\textbf{D wrong:} A summary conclusion is not an evaluation of coverage. \\[3pt]
	\vfill {\tiny Part 2, Set 1}
\end{frame}
% ===== Q48 =====
\begin{frame}[t]{Q48 \hfill {\scriptsize Engagement Planning}}
	\small
	\textbf{Question:} The team plans to reuse last year's work program for a process that was reorganized and moved to a new system. What is the BEST response?
	\begin{enumerate}[label=\Alph*.,itemsep=1pt,topsep=3pt]
			\item Drop the engagement because the process is in flux.
			\item Reuse it but add one extra procedure.
			\item Update the risk assessment for the reorganization, new system, and staff turnover, and tailor the program.
			\item Reuse it unchanged for comparability.
	\end{enumerate}
\end{frame}
\begin{frame}[t]{Q48 -- Explanation}
	\footnotesize
	\textbf{Correct: C.} Significant change alters risk and the controls in operation, and the program has to follow. \\[5pt]
\textbf{A wrong:} Change is a reason to review, not to avoid review. \\[3pt]
\textbf{B wrong:} One added step does not respond to changes in risk. \\[3pt]
\textbf{D wrong:} Comparability does not justify procedures that no longer fit. \\[3pt]
	\vfill {\tiny Part 2, Set 1}
\end{frame}
% ===== Q49 =====
\begin{frame}[t]{Q49 \hfill {\scriptsize Engagement Planning}}
	\small
	\textbf{Question:} Before go-live of a new ERP module, a manager asks internal audit to ``just confirm the controls are fine.'' What is the BEST first response?
	\begin{enumerate}[label=\Alph*.,itemsep=1pt,topsep=3pt]
			\item Accept and decide later whether it is assurance.
			\item Decline all involvement in projects before go-live.
			\item Provide a clean confirmation since the manager asked for it.
			\item Clarify whether the request is for assurance or advisory services and set objectives and scope accordingly, considering any effect on future objectivity.
	\end{enumerate}
\end{frame}
\begin{frame}[t]{Q49 -- Explanation}
	\footnotesize
	\textbf{Correct: D.} The nature of the service determines objectives, independence considerations, and the form of communication. \\[5pt]
\textbf{A wrong:} Deferring the decision leaves expectations and independence unmanaged. \\[3pt]
\textbf{B wrong:} Advisory involvement before go-live can be valuable when scoped properly. \\[3pt]
\textbf{C wrong:} A confirmation without work is an unsupported assurance. \\[3pt]
	\vfill {\tiny Part 2, Set 1}
\end{frame}
% ===== Q50 =====
\begin{frame}[t]{Q50 \hfill {\scriptsize Engagement Planning}}
	\small
	\textbf{Question:} Which engagement objective is BEST written?
	\begin{enumerate}[label=\Alph*.,itemsep=1pt,topsep=3pt]
			\item Evaluate whether controls over third-party access to the customer database prevent and detect unauthorized access during fiscal 2026.
			\item Confirm management's assertion that vendors are secure.
			\item Review third-party access.
			\item Ensure that IT security is adequate.
	\end{enumerate}
\end{frame}
\begin{frame}[t]{Q50 -- Explanation}
	\footnotesize
	\textbf{Correct: A.} A good objective is specific, tied to a risk, and states what will be concluded and over what period. \\[5pt]
\textbf{B wrong:} It starts from management's assertion instead of evaluating the controls. \\[3pt]
\textbf{C wrong:} It is too vague to guide scope or procedures. \\[3pt]
\textbf{D wrong:} It is too broad and uses the word ``ensure,'' which implies a guarantee. \\[3pt]
	\vfill {\tiny Part 2, Set 1}
\end{frame}

%%%%%%%%%% Information Gathering, Analysis and Evaluation %%%%%%%%%%
% ===== Q51 =====
\begin{frame}[t]{Q51 \hfill {\scriptsize Information Gathering, Analysis and Evaluation}}
	\small
	\textbf{Question:} Which evidence is generally the MOST reliable for supporting the existence of customer receivables?
	\begin{enumerate}[label=\Alph*.,itemsep=1pt,topsep=3pt]
			\item An e-mail from the accounts receivable manager.
			\item A confirmation received directly by the auditor from the customer.
			\item An aged trial balance generated by the accounting system.
			\item A verbal explanation from the credit manager.
	\end{enumerate}
\end{frame}
\begin{frame}[t]{Q51 -- Explanation}
	\footnotesize
	\textbf{Correct: B.} Evidence obtained directly from an independent external source is more reliable than internally generated or verbal evidence. \\[5pt]
\textbf{A wrong:} It is internal, written by an interested party, and uncorroborated. \\[3pt]
\textbf{C wrong:} It is internal and depends on the controls over the system. \\[3pt]
\textbf{D wrong:} Verbal statements are the least reliable form. \\[3pt]
	\vfill {\tiny Part 2, Set 1}
\end{frame}
% ===== Q52 =====
\begin{frame}[t]{Q52 \hfill {\scriptsize Information Gathering, Analysis and Evaluation}}
	\small
	\textbf{Question:} The auditor suspects ``ghost employees.'' Which procedure BEST addresses this?
	\begin{enumerate}[label=\Alph*.,itemsep=1pt,topsep=3pt]
			\item Recalculate gross pay for a sample of employees.
			\item Select entries from the payroll register and trace them to HR records, approvals, and evidence of actual employment or payment.
			\item Compare total payroll cost to budget.
			\item Select from the active HR list and trace each person to the payroll register.
	\end{enumerate}
\end{frame}
\begin{frame}[t]{Q52 -- Explanation}
	\footnotesize
	\textbf{Correct: B.} Existence of fictitious payees is tested by working from what was paid back to supporting records. \\[5pt]
\textbf{A wrong:} Recalculation tests accuracy, not whether the person exists. \\[3pt]
\textbf{C wrong:} Budget comparison cannot reveal individual fictitious entries. \\[3pt]
\textbf{D wrong:} That direction tests completeness, such as unpaid employees, not fictitious payees. \\[3pt]
	\vfill {\tiny Part 2, Set 1}
\end{frame}
% ===== Q53 =====
\begin{frame}[t]{Q53 \hfill {\scriptsize Information Gathering, Analysis and Evaluation}}
	\small
	\textbf{Question:} The auditor wants evidence that all fixed asset additions have been recorded. Which approach is BEST?
	\begin{enumerate}[label=\Alph*.,itemsep=1pt,topsep=3pt]
			\item Select from the asset register and trace to invoices.
			\item Ask the controller whether all additions were recorded.
			\item Compare depreciation expense to the budget.
			\item Select from receiving reports and vendor invoices for capital items and trace them to the asset register.
	\end{enumerate}
\end{frame}
\begin{frame}[t]{Q53 -- Explanation}
	\footnotesize
	\textbf{Correct: D.} Completeness is tested by starting from independent source documents and tracing forward to the records. \\[5pt]
\textbf{A wrong:} That tests existence of recorded items, not unrecorded ones. \\[3pt]
\textbf{B wrong:} Inquiry alone is weak evidence. \\[3pt]
\textbf{C wrong:} Depreciation versus budget is only a general analytical indicator. \\[3pt]
	\vfill {\tiny Part 2, Set 1}
\end{frame}
% ===== Q54 =====
\begin{frame}[t]{Q54 \hfill {\scriptsize Information Gathering, Analysis and Evaluation}}
	\small
	\textbf{Question:} Which procedure is MOST effective in identifying unrecorded liabilities at year-end?
	\begin{enumerate}[label=\Alph*.,itemsep=1pt,topsep=3pt]
			\item Confirm balances with the largest suppliers only.
			\item Compare accounts payable to last year's balance.
			\item Examine disbursements and receiving reports after year-end and trace those relating to prior-period goods or services to recorded payables.
			\item Vouch recorded payables to supporting invoices.
	\end{enumerate}
\end{frame}
\begin{frame}[t]{Q54 -- Explanation}
	\footnotesize
	\textbf{Correct: C.} A search for unrecorded items works from subsequent events back to the period under review. \\[5pt]
\textbf{A wrong:} Confirmations of large suppliers miss unrecorded amounts, including small suppliers. \\[3pt]
\textbf{B wrong:} A prior-year comparison is too coarse to reveal omissions. \\[3pt]
\textbf{D wrong:} Vouching tests items already recorded and cannot find omissions. \\[3pt]
	\vfill {\tiny Part 2, Set 1}
\end{frame}
% ===== Q55 =====
\begin{frame}[t]{Q55 \hfill {\scriptsize Information Gathering, Analysis and Evaluation}}
	\small
	\textbf{Question:} A customer does not respond to a receivable confirmation request. Which alternative procedure is MOST appropriate?
	\begin{enumerate}[label=\Alph*.,itemsep=1pt,topsep=3pt]
			\item Accept the sales manager's assurance that the customer is reliable.
			\item Examine subsequent cash receipts and shipping or service documents supporting the sale.
			\item Treat non-response as confirmation of the balance.
			\item Remove the item from the sample.
	\end{enumerate}
\end{frame}
\begin{frame}[t]{Q55 -- Explanation}
	\footnotesize
	\textbf{Correct: B.} Alternatives should provide evidence that the receivable existed and was valid. \\[5pt]
\textbf{A wrong:} An assurance from a party with an interest is weak. \\[3pt]
\textbf{C wrong:} Silence is not evidence of agreement. \\[3pt]
\textbf{D wrong:} Dropping non-responses biases the sample toward easy items. \\[3pt]
	\vfill {\tiny Part 2, Set 1}
\end{frame}
% ===== Q56 =====
\begin{frame}[t]{Q56 \hfill {\scriptsize Information Gathering, Analysis and Evaluation}}
	\small
	\textbf{Question:} An automated three-way match has a 2\% price tolerance. Which test BEST shows it is working?
	\begin{enumerate}[label=\Alph*.,itemsep=1pt,topsep=3pt]
			\item Check that the match report was generated on schedule.
			\item Review the policy that describes the tolerance.
			\item Submit or observe test transactions above and below the tolerance, or inspect configuration and its change history, together with ITGC evidence.
			\item Ask the AP manager whether exceptions occur.
	\end{enumerate}
\end{frame}
\begin{frame}[t]{Q56 -- Explanation}
	\footnotesize
	\textbf{Correct: C.} A configured control is tested by seeing how the system treats conditions it should accept and reject. \\[5pt]
\textbf{A wrong:} A report running on time does not show it applies the tolerance. \\[3pt]
\textbf{B wrong:} A policy shows intent, not that the system enforces it. \\[3pt]
\textbf{D wrong:} Inquiry offers little evidence on system behavior. \\[3pt]
	\vfill {\tiny Part 2, Set 1}
\end{frame}
% ===== Q57 =====
\begin{frame}[t]{Q57 \hfill {\scriptsize Information Gathering, Analysis and Evaluation}}
	\small
	\textbf{Question:} Which procedure BEST tests segregation of duties in an ERP system?
	\begin{enumerate}[label=\Alph*.,itemsep=1pt,topsep=3pt]
			\item Review the organization chart.
			\item Ask supervisors whether staff stay within their duties.
			\item Observe a single transaction from start to finish.
			\item Obtain an extract of roles and access directly from the system, compare it to a conflict matrix, and check whether conflicts were used in transactions.
	\end{enumerate}
\end{frame}
\begin{frame}[t]{Q57 -- Explanation}
	\footnotesize
	\textbf{Correct: D.} System access determines what people can actually do, regardless of titles. \\[5pt]
\textbf{A wrong:} An org chart shows reporting lines, not system access. \\[3pt]
\textbf{B wrong:} Supervisors' views do not reveal technical access rights. \\[3pt]
\textbf{C wrong:} One observation cannot show what access exists across users. \\[3pt]
	\vfill {\tiny Part 2, Set 1}
\end{frame}
% ===== Q58 =====
\begin{frame}[t]{Q58 \hfill {\scriptsize Information Gathering, Analysis and Evaluation}}
	\small
	\textbf{Question:} To test whether terminated employees retain system access, which procedure is BEST?
	\begin{enumerate}[label=\Alph*.,itemsep=1pt,topsep=3pt]
			\item Test access for employees currently on the HR active list.
			\item Compare a termination list from HR with a system-generated list of active accounts and review activity after termination dates.
			\item Ask IT whether access is removed promptly.
			\item Review the policy requiring prompt removal.
	\end{enumerate}
\end{frame}
\begin{frame}[t]{Q58 -- Explanation}
	\footnotesize
	\textbf{Correct: B.} The control is tested by comparing independent sources and looking for use after separation. \\[5pt]
\textbf{A wrong:} That tests the wrong population. \\[3pt]
\textbf{C wrong:} Inquiry alone is not sufficient evidence. \\[3pt]
\textbf{D wrong:} The policy shows design, not operation. \\[3pt]
	\vfill {\tiny Part 2, Set 1}
\end{frame}
% ===== Q59 =====
\begin{frame}[t]{Q59 \hfill {\scriptsize Information Gathering, Analysis and Evaluation}}
	\small
	\textbf{Question:} Which journal-entry characteristic MOST strongly indicates possible management override?
	\begin{enumerate}[label=\Alph*.,itemsep=1pt,topsep=3pt]
			\item Routine invoice entries posted by accounts payable clerks.
			\item Entries with sequential numbers in a standard format.
			\item Manual entries posted by a senior finance executive after period end to revenue and reserve accounts with no description.
			\item Recurring automated depreciation entries.
	\end{enumerate}
\end{frame}
\begin{frame}[t]{Q59 -- Explanation}
	\footnotesize
	\textbf{Correct: C.} Override typically involves unusual users, timing, accounts, and weak support. \\[5pt]
\textbf{A wrong:} Routine entries by process staff are expected. \\[3pt]
\textbf{B wrong:} Sequential numbering is normal and not a risk indicator. \\[3pt]
\textbf{D wrong:} Automated recurring entries are normal and rule-based. \\[3pt]
	\vfill {\tiny Part 2, Set 1}
\end{frame}
% ===== Q60 =====
\begin{frame}[t]{Q60 \hfill {\scriptsize Information Gathering, Analysis and Evaluation}}
	\small
	\textbf{Question:} To test a bank reconciliation control, which procedure is BEST?
	\begin{enumerate}[label=\Alph*.,itemsep=1pt,topsep=3pt]
			\item Reperform the reconciliation from bank and ledger data, and investigate aging reconciling items and review evidence.
			\item Confirm that the reconciliation policy exists.
			\item Compare the cash balance to last year.
			\item Ask the preparer whether differences are resolved.
	\end{enumerate}
\end{frame}
\begin{frame}[t]{Q60 -- Explanation}
	\footnotesize
	\textbf{Correct: A.} Reperformance shows whether the reconciliation works and whether differences are investigated. \\[5pt]
\textbf{B wrong:} A policy review tests design only. \\[3pt]
\textbf{C wrong:} A year-over-year balance comparison does not test the control. \\[3pt]
\textbf{D wrong:} The preparer's statement is not independent evidence. \\[3pt]
	\vfill {\tiny Part 2, Set 1}
\end{frame}
% ===== Q61 =====
\begin{frame}[t]{Q61 \hfill {\scriptsize Information Gathering, Analysis and Evaluation}}
	\small
	\textbf{Question:} Which procedure BEST tests the completeness of the IT change-management log?
	\begin{enumerate}[label=\Alph*.,itemsep=1pt,topsep=3pt]
			\item Ask the IT director whether all changes are logged.
			\item Review the change-management policy.
			\item Count the tickets in the ticketing system.
			\item Compare changes identified from independent source such as deployment tools or version control to approved change tickets.
	\end{enumerate}
\end{frame}
\begin{frame}[t]{Q61 -- Explanation}
	\footnotesize
	\textbf{Correct: D.} To find unlogged changes, the population must come from a source independent of the log being tested. \\[5pt]
\textbf{A wrong:} Inquiry alone cannot detect unlogged changes. \\[3pt]
\textbf{B wrong:} A policy shows design, not whether changes bypassed the process. \\[3pt]
\textbf{C wrong:} Counting tickets assumes the ticket population is complete. \\[3pt]
	\vfill {\tiny Part 2, Set 1}
\end{frame}
% ===== Q62 =====
\begin{frame}[t]{Q62 \hfill {\scriptsize Information Gathering, Analysis and Evaluation}}
	\small
	\textbf{Question:} Before running analytics on data extracted from the ERP, what should the auditor do?
	\begin{enumerate}[label=\Alph*.,itemsep=1pt,topsep=3pt]
			\item Ask the process owner to confirm the data is correct.
			\item Assume system-generated data is reliable.
			\item Run the analysis first and validate only if anomalies appear.
			\item Validate completeness and accuracy of the extract, such as reconciling record counts and control totals to source.
	\end{enumerate}
\end{frame}
\begin{frame}[t]{Q62 -- Explanation}
	\footnotesize
	\textbf{Correct: D.} Conclusions from analytics depend on the data they are based on. \\[5pt]
\textbf{A wrong:} A process owner's statement is not independent validation. \\[3pt]
\textbf{B wrong:} System generation does not guarantee completeness of an extract. \\[3pt]
\textbf{C wrong:} Validating after the fact risks wrong conclusions in the meantime. \\[3pt]
	\vfill {\tiny Part 2, Set 1}
\end{frame}
% ===== Q63 =====
\begin{frame}[t]{Q63 \hfill {\scriptsize Information Gathering, Analysis and Evaluation}}
	\small
	\textbf{Question:} Population for a test of new-vendor approvals is a log kept by the process owner in a spreadsheet. What should the auditor do before sampling?
	\begin{enumerate}[label=\Alph*.,itemsep=1pt,topsep=3pt]
			\item Reconcile the log to an independent source, such as the vendor master-file creation history.
			\item Sample from the log since it is the only list available.
			\item Ask the owner to certify the log is complete.
			\item Increase the sample size to compensate for uncertainty.
	\end{enumerate}
\end{frame}
\begin{frame}[t]{Q63 -- Explanation}
	\footnotesize
	\textbf{Correct: A.} A sample is only as good as the population it is drawn from. \\[5pt]
\textbf{B wrong:} If the log omits vendors, the sample cannot detect it. \\[3pt]
\textbf{C wrong:} A certification from the party being tested is not independent. \\[3pt]
\textbf{D wrong:} A bigger sample from an incomplete list is still incomplete. \\[3pt]
	\vfill {\tiny Part 2, Set 1}
\end{frame}
% ===== Q64 =====
\begin{frame}[t]{Q64 \hfill {\scriptsize Information Gathering, Analysis and Evaluation}}
	\small
	\textbf{Question:} Which procedure BEST tests the integrity of a spreadsheet used in a critical financial estimate?
	\begin{enumerate}[label=\Alph*.,itemsep=1pt,topsep=3pt]
			\item Check the last-modified date.
			\item Independently recalculate results and review formula protection, version history, and reconciliation of inputs to source.
			\item Obtain the preparer's statement that the formulas are correct.
			\item Re-add the totals only.
	\end{enumerate}
\end{frame}
\begin{frame}[t]{Q64 -- Explanation}
	\footnotesize
	\textbf{Correct: B.} Integrity depends on formulas, inputs, and change control, not just arithmetic. \\[5pt]
\textbf{A wrong:} A date says nothing about correctness. \\[3pt]
\textbf{C wrong:} A statement from the preparer is not independent evidence. \\[3pt]
\textbf{D wrong:} Re-adding totals does not test formulas or inputs. \\[3pt]
	\vfill {\tiny Part 2, Set 1}
\end{frame}
% ===== Q65 =====
\begin{frame}[t]{Q65 \hfill {\scriptsize Information Gathering, Analysis and Evaluation}}
	\small
	\textbf{Question:} Which dataset is LEAST suitable for a Benford's Law test?
	\begin{enumerate}[label=\Alph*.,itemsep=1pt,topsep=3pt]
			\item Transaction amounts across many accounts.
			\item Employee numbers assigned sequentially.
			\item Expense claim amounts.
			\item Vendor invoice amounts across many suppliers.
	\end{enumerate}
\end{frame}
\begin{frame}[t]{Q65 -- Explanation}
	\footnotesize
	\textbf{Correct: B.} Benford's Law applies to naturally occurring numbers across several orders of magnitude, not assigned sequences. \\[5pt]
\textbf{A wrong:} Transaction amounts across accounts are typically suitable. \\[3pt]
\textbf{C wrong:} Expense amounts commonly fit the pattern. \\[3pt]
\textbf{D wrong:} Invoice amounts are a typical suitable dataset. \\[3pt]
	\vfill {\tiny Part 2, Set 1}
\end{frame}
% ===== Q66 =====
\begin{frame}[t]{Q66 \hfill {\scriptsize Information Gathering, Analysis and Evaluation}}
	\small
	\textbf{Question:} An anomaly-detection model flags 2,000 payments for review. What is the BEST way to proceed?
	\begin{enumerate}[label=\Alph*.,itemsep=1pt,topsep=3pt]
			\item Use risk scoring to prioritize, investigate a defined set, document model logic and thresholds, and treat flags as leads, not findings.
			\item Review only the five highest-scoring items.
			\item Discard flags that the process owner says are normal.
			\item Treat each flagged payment as a finding.
	\end{enumerate}
\end{frame}
\begin{frame}[t]{Q66 -- Explanation}
	\footnotesize
	\textbf{Correct: A.} Flags indicate where to look, and conclusions require corroborating evidence. \\[5pt]
\textbf{B wrong:} Five items are too few to evaluate the model or the population. \\[3pt]
\textbf{C wrong:} Dismissal based on owner statements removes independence. \\[3pt]
\textbf{D wrong:} A flag alone is not evidence of a control failure or error. \\[3pt]
	\vfill {\tiny Part 2, Set 1}
\end{frame}
% ===== Q67 =====
\begin{frame}[t]{Q67 \hfill {\scriptsize Information Gathering, Analysis and Evaluation}}
	\small
	\textbf{Question:} An auditor wants to use a generative AI tool to summarize interviews and draft test conclusions. Which approach is BEST?
	\begin{enumerate}[label=\Alph*.,itemsep=1pt,topsep=3pt]
			\item Upload workpapers to a public tool to save time.
			\item Validate outputs against source evidence, protect confidential data, and document the use and review.
			\item Treat the tool's output as audit evidence.
			\item Use AI output in place of supervisory review.
	\end{enumerate}
\end{frame}
\begin{frame}[t]{Q67 -- Explanation}
	\footnotesize
	\textbf{Correct: B.} AI can support the work but cannot replace evidence, confidentiality, or review. \\[5pt]
\textbf{A wrong:} Public tools can expose confidential engagement information. \\[3pt]
\textbf{C wrong:} Output is a derived product, not evidence. \\[3pt]
\textbf{D wrong:} Supervisory review is still required for the work. \\[3pt]
	\vfill {\tiny Part 2, Set 1}
\end{frame}
% ===== Q68 =====
\begin{frame}[t]{Q68 \hfill {\scriptsize Information Gathering, Analysis and Evaluation}}
	\small
	\textbf{Question:} An embedded continuous-monitoring routine alerts on high-risk payments. What is the MOST important consideration before relying on it?
	\begin{enumerate}[label=\Alph*.,itemsep=1pt,topsep=3pt]
			\item Confirm the dashboard looks professional.
			\item Validate its logic, thresholds, data coverage, and the controls over changes to the routine.
			\item Confirm the vendor has other clients.
			\item Confirm alerts are sent by e-mail.
	\end{enumerate}
\end{frame}
\begin{frame}[t]{Q68 -- Explanation}
	\footnotesize
	\textbf{Correct: B.} Automated monitoring misses risks when logic or coverage is flawed, so it is validated like any control. \\[5pt]
\textbf{A wrong:} Presentation does not show effectiveness. \\[3pt]
\textbf{C wrong:} Other clients do not show the logic works here. \\[3pt]
\textbf{D wrong:} The delivery channel is irrelevant to detection quality. \\[3pt]
	\vfill {\tiny Part 2, Set 1}
\end{frame}
% ===== Q69 =====
\begin{frame}[t]{Q69 \hfill {\scriptsize Information Gathering, Analysis and Evaluation}}
	\small
	\textbf{Question:} A RACI chart for an approval step shows two roles marked ``Accountable.'' What does this MOST likely indicate?
	\begin{enumerate}[label=\Alph*.,itemsep=1pt,topsep=3pt]
			\item Appropriate segregation of duties.
			\item Strong redundancy that improves control.
			\item A documentation error with no control effect.
			\item Unclear accountability, a potential control-design weakness.
	\end{enumerate}
\end{frame}
\begin{frame}[t]{Q69 -- Explanation}
	\footnotesize
	\textbf{Correct: D.} A single point of accountability is expected, and two can mean each assumes the other will act. \\[5pt]
\textbf{A wrong:} Segregation is about separating incompatible duties, not duplicating accountability. \\[3pt]
\textbf{B wrong:} Shared accountability often leads to gaps rather than redundancy. \\[3pt]
\textbf{C wrong:} It can reflect real confusion about who owns the outcome. \\[3pt]
	\vfill {\tiny Part 2, Set 1}
\end{frame}
% ===== Q70 =====
\begin{frame}[t]{Q70 \hfill {\scriptsize Information Gathering, Analysis and Evaluation}}
	\small
	\textbf{Question:} The auditor needs to find nonstandard termination clauses across 3,000 PDF contracts. Which technique is MOST suitable?
	\begin{enumerate}[label=\Alph*.,itemsep=1pt,topsep=3pt]
			\item Regression on invoice amounts.
			\item Benford's Law on contract values.
			\item Text analytics on the unstructured contract data, followed by review of flagged documents.
			\item Stratified sampling on contract value only.
	\end{enumerate}
\end{frame}
\begin{frame}[t]{Q70 -- Explanation}
	\footnotesize
	\textbf{Correct: C.} Clause content is unstructured text and requires text-based methods. \\[5pt]
\textbf{A wrong:} Invoice regression does not examine contract text. \\[3pt]
\textbf{B wrong:} Numeric digit patterns do not reveal clause wording. \\[3pt]
\textbf{D wrong:} Value-based sampling would miss clause variations in lower-value contracts. \\[3pt]
	\vfill {\tiny Part 2, Set 1}
\end{frame}
% ===== Q71 =====
\begin{frame}[t]{Q71 \hfill {\scriptsize Information Gathering, Analysis and Evaluation}}
	\small
	\textbf{Question:} A team builds a model estimating the likelihood that new vendors will submit duplicate invoices. This is an example of which analytics type?
	\begin{enumerate}[label=\Alph*.,itemsep=1pt,topsep=3pt]
			\item Diagnostic analytics.
			\item Predictive analytics.
			\item Prescriptive analytics.
			\item Descriptive analytics.
	\end{enumerate}
\end{frame}
\begin{frame}[t]{Q71 -- Explanation}
	\footnotesize
	\textbf{Correct: B.} Predictive analytics estimates what is likely to happen. \\[5pt]
\textbf{A wrong:} Diagnostic analytics explains why something happened. \\[3pt]
\textbf{C wrong:} Prescriptive analytics recommends actions to take. \\[3pt]
\textbf{D wrong:} Descriptive analytics summarizes what happened. \\[3pt]
	\vfill {\tiny Part 2, Set 1}
\end{frame}
% ===== Q72 =====
\begin{frame}[t]{Q72 \hfill {\scriptsize Information Gathering, Analysis and Evaluation}}
	\small
	\textbf{Question:} Gross margin fell 4 points. Which analytical review approach is BEST?
	\begin{enumerate}[label=\Alph*.,itemsep=1pt,topsep=3pt]
			\item Build an independent expectation using drivers such as volume, price, and mix, compare it with results, and investigate differences above a defined threshold.
			\item Accept management's explanation as sufficient evidence.
			\item Compare with last year only.
			\item Ignore variances below last year's threshold.
	\end{enumerate}
\end{frame}
\begin{frame}[t]{Q72 -- Explanation}
	\footnotesize
	\textbf{Correct: A.} Analytics work when the expectation is independent and variances are investigated against corroborating evidence. \\[5pt]
\textbf{B wrong:} An explanation is a lead that needs corroboration. \\[3pt]
\textbf{C wrong:} A prior-year comparison does not reflect current drivers. \\[3pt]
\textbf{D wrong:} Thresholds should reflect the current risk and tolerable levels, not old ones. \\[3pt]
	\vfill {\tiny Part 2, Set 1}
\end{frame}
% ===== Q73 =====
\begin{frame}[t]{Q73 \hfill {\scriptsize Information Gathering, Analysis and Evaluation}}
	\small
	\textbf{Question:} Which statement is a ROOT CAUSE in a finding about unauthorized overrides?
	\begin{enumerate}[label=\Alph*.,itemsep=1pt,topsep=3pt]
			\item Overrides occurred in 18\% of transactions.
			\item The organization risks regulatory fines.
			\item The policy requires approval for overrides.
			\item Overrides were possible because the system had no mandatory reason code and no one reviewed override reports.
	\end{enumerate}
\end{frame}
\begin{frame}[t]{Q73 -- Explanation}
	\footnotesize
	\textbf{Correct: D.} A root cause explains why the condition occurred. \\[5pt]
\textbf{A wrong:} This describes the condition. \\[3pt]
\textbf{B wrong:} This describes the effect. \\[3pt]
\textbf{C wrong:} This states the criterion. \\[3pt]
	\vfill {\tiny Part 2, Set 1}
\end{frame}
% ===== Q74 =====
\begin{frame}[t]{Q74 \hfill {\scriptsize Information Gathering, Analysis and Evaluation}}
	\small
	\textbf{Question:} Which approach to evaluating the significance of a finding is BEST?
	\begin{enumerate}[label=\Alph*.,itemsep=1pt,topsep=3pt]
			\item Consider impact, likelihood, and pervasiveness, along with qualitative factors such as legal exposure or fraud indicators.
			\item Rate significance by dollar amount only.
			\item Rate by the number of occurrences.
			\item Rate by management's willingness to fix it.
	\end{enumerate}
\end{frame}
\begin{frame}[t]{Q74 -- Explanation}
	\footnotesize
	\textbf{Correct: A.} Significance reflects both quantitative and qualitative factors. \\[5pt]
\textbf{B wrong:} Small amounts can still indicate serious failures. \\[3pt]
\textbf{C wrong:} Frequency alone ignores impact. \\[3pt]
\textbf{D wrong:} Willingness to fix does not change how serious the issue is. \\[3pt]
	\vfill {\tiny Part 2, Set 1}
\end{frame}
% ===== Q75 =====
\begin{frame}[t]{Q75 \hfill {\scriptsize Information Gathering, Analysis and Evaluation}}
	\small
	\textbf{Question:} A judgmental sample of the 25 largest items showed no exceptions. Which conclusion is supportable?
	\begin{enumerate}[label=\Alph*.,itemsep=1pt,topsep=3pt]
			\item The population error rate is below 4\%.
			\item The result is supported at a 95\% confidence level.
			\item The control is effective for all items.
			\item Only that those 25 items had no exceptions; the result cannot be projected to the population.
	\end{enumerate}
\end{frame}
\begin{frame}[t]{Q75 -- Explanation}
	\footnotesize
	\textbf{Correct: D.} Judgmental samples describe the items selected only. \\[5pt]
\textbf{A wrong:} No statistical basis exists for the rate. \\[3pt]
\textbf{B wrong:} Confidence levels require a statistical design. \\[3pt]
\textbf{C wrong:} Effectiveness across the population is not shown by a biased selection. \\[3pt]
	\vfill {\tiny Part 2, Set 1}
\end{frame}
% ===== Q76 =====
\begin{frame}[t]{Q76 \hfill {\scriptsize Information Gathering, Analysis and Evaluation}}
	\small
	\textbf{Question:} Which description matches ``sufficient'' information under the Standards?
	\begin{enumerate}[label=\Alph*.,itemsep=1pt,topsep=3pt]
			\item Information consistent with the engagement's objectives and scope.
			\item Information that is factual, adequate, and convincing so a prudent, informed person would reach the same conclusions.
			\item Information that helps the organization meet its goals.
			\item Information that is the best attainable using appropriate engagement techniques.
	\end{enumerate}
\end{frame}
\begin{frame}[t]{Q76 -- Explanation}
	\footnotesize
	\textbf{Correct: B.} Sufficiency concerns whether the information is enough to be convincing. \\[5pt]
\textbf{A wrong:} This describes relevance. \\[3pt]
\textbf{C wrong:} This describes usefulness. \\[3pt]
\textbf{D wrong:} This describes reliability. \\[3pt]
	\vfill {\tiny Part 2, Set 1}
\end{frame}
% ===== Q77 =====
\begin{frame}[t]{Q77 \hfill {\scriptsize Information Gathering, Analysis and Evaluation}}
	\small
	\textbf{Question:} An auditor leaving the organization wants to keep copies of workpapers for personal reference after removing client names. What is the correct response?
	\begin{enumerate}[label=\Alph*.,itemsep=1pt,topsep=3pt]
			\item Not permitted; workpapers belong to the organization and access and retention are governed by policy and law.
			\item Permitted for the auditor's own work only.
			\item Permitted if the names are removed.
			\item Permitted for 12 months after departure.
	\end{enumerate}
\end{frame}
\begin{frame}[t]{Q77 -- Explanation}
	\footnotesize
	\textbf{Correct: A.} Engagement records are organizational property whose access is controlled. \\[5pt]
\textbf{B wrong:} Ownership does not depend on who prepared the file. \\[3pt]
\textbf{C wrong:} Removing names does not change ownership or confidentiality. \\[3pt]
\textbf{D wrong:} No time-limited personal right exists. \\[3pt]
	\vfill {\tiny Part 2, Set 1}
\end{frame}
% ===== Q78 =====
\begin{frame}[t]{Q78 \hfill {\scriptsize Information Gathering, Analysis and Evaluation}}
	\small
	\textbf{Question:} While testing, an auditor finds evidence that the CFO directed a reclassification to meet a loan covenant. What is the MOST appropriate action?
	\begin{enumerate}[label=\Alph*.,itemsep=1pt,topsep=3pt]
			\item Finish the engagement and mention it in the final report.
			\item Raise skepticism, expand testing as needed, and promptly inform the CAE under the fraud-escalation protocol, without alerting implicated management.
			\item Discuss it with the CFO first to get an explanation.
			\item Narrow the scope to avoid the area.
	\end{enumerate}
\end{frame}
\begin{frame}[t]{Q78 -- Explanation}
	\footnotesize
	\textbf{Correct: B.} Possible fraud indicators involving management require careful escalation, not discussion with the implicated party. \\[5pt]
\textbf{A wrong:} Delay may harm the organization. \\[3pt]
\textbf{C wrong:} Tipping off the CFO risks evidence loss and interference. \\[3pt]
\textbf{D wrong:} Narrowing scope leaves the issue unexamined. \\[3pt]
	\vfill {\tiny Part 2, Set 1}
\end{frame}
% ===== Q79 =====
\begin{frame}[t]{Q79 \hfill {\scriptsize Information Gathering, Analysis and Evaluation}}
	\small
	\textbf{Question:} An auditor requests system logs and learns they were deleted the day before. What is the BEST response?
	\begin{enumerate}[label=\Alph*.,itemsep=1pt,topsep=3pt]
			\item Confront the staff aggressively.
			\item Ask staff to recreate the logs from memory.
			\item Document the circumstances, preserve remaining evidence, and notify the CAE so the matter can be escalated to appropriate parties such as legal.
			\item Accept the explanation and continue as planned.
	\end{enumerate}
\end{frame}
\begin{frame}[t]{Q79 -- Explanation}
	\footnotesize
	\textbf{Correct: C.} Loss of evidence may indicate wrongdoing or a legal issue and requires preservation and escalation. \\[5pt]
\textbf{A wrong:} Confrontation can compromise evidence and any later inquiry. \\[3pt]
\textbf{B wrong:} Recreated records are not reliable evidence. \\[3pt]
\textbf{D wrong:} Continuing without escalation ignores potential misconduct. \\[3pt]
	\vfill {\tiny Part 2, Set 1}
\end{frame}
% ===== Q80 =====
\begin{frame}[t]{Q80 \hfill {\scriptsize Information Gathering, Analysis and Evaluation}}
	\small
	\textbf{Question:} The process owner says reviews happen monthly, but reviewer e-mails show quarterly reviews. What is the BEST action?
	\begin{enumerate}[label=\Alph*.,itemsep=1pt,topsep=3pt]
			\item Obtain system timestamps or other independent evidence, reperform where possible, and document how the conflict was resolved.
			\item Report both versions without further work.
			\item Accept the owner's statement since it is the more recent.
			\item Accept the e-mails and ignore the statement.
	\end{enumerate}
\end{frame}
\begin{frame}[t]{Q80 -- Explanation}
	\footnotesize
	\textbf{Correct: A.} Conflicting evidence must be resolved through more reliable sources. \\[5pt]
\textbf{B wrong:} Presenting unresolved conflict leaves the issue unevaluated. \\[3pt]
\textbf{C wrong:} Recency does not make a statement more reliable. \\[3pt]
\textbf{D wrong:} One source cannot be chosen without testing. \\[3pt]
	\vfill {\tiny Part 2, Set 1}
\end{frame}
% ===== Q81 =====
\begin{frame}[t]{Q81 \hfill {\scriptsize Information Gathering, Analysis and Evaluation}}
	\small
	\textbf{Question:} A control deficiency is found, but management points to a compensating control. What should the auditor do?
	\begin{enumerate}[label=\Alph*.,itemsep=1pt,topsep=3pt]
			\item Ignore the deficiency because compensation exists.
			\item Test whether the compensating control is precise enough and operating effectively, then evaluate the remaining deficiency.
			\item Report the deficiency without considering the compensating control.
			\item Assume the compensating control works.
	\end{enumerate}
\end{frame}
\begin{frame}[t]{Q81 -- Explanation}
	\footnotesize
	\textbf{Correct: B.} A compensating control reduces the risk only if it actually works. \\[5pt]
\textbf{A wrong:} The deficiency may remain significant if the compensating control is weak. \\[3pt]
\textbf{C wrong:} Omitting it may overstate the exposure. \\[3pt]
\textbf{D wrong:} Management's claim needs evidence. \\[3pt]
	\vfill {\tiny Part 2, Set 1}
\end{frame}
% ===== Q82 =====
\begin{frame}[t]{Q82 \hfill {\scriptsize Information Gathering, Analysis and Evaluation}}
	\small
	\textbf{Question:} An engagement tests compliance with a debt covenant requiring interest cover of at least 3.0. Which evidence is MOST appropriate?
	\begin{enumerate}[label=\Alph*.,itemsep=1pt,topsep=3pt]
			\item Recompute the ratio from general ledger data using the definitions in the loan agreement.
			\item Review management's compliance certificate.
			\item Confirm verbally with the treasurer.
			\item Compare the ratio to industry averages.
	\end{enumerate}
\end{frame}
\begin{frame}[t]{Q82 -- Explanation}
	\footnotesize
	\textbf{Correct: A.} Compliance depends on the agreement's definitions, so recomputation from source is strongest. \\[5pt]
\textbf{B wrong:} The certificate is management's assertion and may use different definitions. \\[3pt]
\textbf{C wrong:} A verbal confirmation is weak evidence. \\[3pt]
\textbf{D wrong:} Industry averages do not show compliance with the agreement. \\[3pt]
	\vfill {\tiny Part 2, Set 1}
\end{frame}
% ===== Q83 =====
\begin{frame}[t]{Q83 \hfill {\scriptsize Information Gathering, Analysis and Evaluation}}
	\small
	\textbf{Question:} Analytics show 1.2\% of invoices were paid without a purchase order. What is the BEST next step?
	\begin{enumerate}[label=\Alph*.,itemsep=1pt,topsep=3pt]
			\item Increase the audit budget.
			\item Conclude it is acceptable because it is under 5\%.
			\item Treat all as fraud.
			\item Categorize the exceptions by cause and test whether each was authorized.
	\end{enumerate}
\end{frame}
\begin{frame}[t]{Q83 -- Explanation}
	\footnotesize
	\textbf{Correct: D.} Understanding why exceptions arise and whether they were approved determines significance. \\[5pt]
\textbf{A wrong:} Budget changes do not resolve the question. \\[3pt]
\textbf{B wrong:} Arbitrary thresholds do not show exceptions were authorized. \\[3pt]
\textbf{C wrong:} Exceptions may be legitimate, such as emergency purchases. \\[3pt]
	\vfill {\tiny Part 2, Set 1}
\end{frame}
% ===== Q84 =====
\begin{frame}[t]{Q84 \hfill {\scriptsize Information Gathering, Analysis and Evaluation}}
	\small
	\textbf{Question:} A vulnerability scan lists 400 high-severity vulnerabilities. How should the auditor evaluate significance?
	\begin{enumerate}[label=\Alph*.,itemsep=1pt,topsep=3pt]
			\item Count them and report the total as the finding.
			\item Ignore them because scans have false positives.
			\item Consider asset criticality, exposure, exploitability, and compensating controls.
			\item Consider only those on non-production servers.
	\end{enumerate}
\end{frame}
\begin{frame}[t]{Q84 -- Explanation}
	\footnotesize
	\textbf{Correct: C.} Significance depends on context rather than raw counts. \\[5pt]
\textbf{A wrong:} A count does not show which are material. \\[3pt]
\textbf{B wrong:} False positives are filtered, not ignored. \\[3pt]
\textbf{D wrong:} Non-production systems are not the main concern. \\[3pt]
	\vfill {\tiny Part 2, Set 1}
\end{frame}
% ===== Q85 =====
\begin{frame}[t]{Q85 \hfill {\scriptsize Information Gathering, Analysis and Evaluation}}
	\small
	\textbf{Question:} To test customer data deletion under a retention schedule, which procedure is BEST?
	\begin{enumerate}[label=\Alph*.,itemsep=1pt,topsep=3pt]
			\item Ask the privacy officer whether deletion occurs.
			\item Review a sample of records still within retention.
			\item Select records past their retention period from the system and verify deletion, including backups and third-party copies where applicable.
			\item Confirm the retention schedule exists.
	\end{enumerate}
\end{frame}
\begin{frame}[t]{Q85 -- Explanation}
	\footnotesize
	\textbf{Correct: C.} Testing deletion requires looking at records that should no longer exist. \\[5pt]
\textbf{A wrong:} Inquiry is weak evidence. \\[3pt]
\textbf{B wrong:} Records within retention cannot show failure to delete. \\[3pt]
\textbf{D wrong:} A schedule shows design only. \\[3pt]
	\vfill {\tiny Part 2, Set 1}
\end{frame}
% ===== Q86 =====
\begin{frame}[t]{Q86 \hfill {\scriptsize Information Gathering, Analysis and Evaluation}}
	\small
	\textbf{Question:} Which evidence BEST shows that backups are restorable?
	\begin{enumerate}[label=\Alph*.,itemsep=1pt,topsep=3pt]
			\item Results of an actual restoration test, with restored data compared with source for integrity.
			\item The backup schedule.
			\item Backup job completion logs.
			\item A contract with an offsite storage vendor.
	\end{enumerate}
\end{frame}
\begin{frame}[t]{Q86 -- Explanation}
	\footnotesize
	\textbf{Correct: A.} Only an actual restoration demonstrates that data can be recovered. \\[5pt]
\textbf{B wrong:} A schedule shows intent only. \\[3pt]
\textbf{C wrong:} Completion means the job ran, not that data can be restored. \\[3pt]
\textbf{D wrong:} A contract shows storage, not recoverability. \\[3pt]
	\vfill {\tiny Part 2, Set 1}
\end{frame}
% ===== Q87 =====
\begin{frame}[t]{Q87 \hfill {\scriptsize Information Gathering, Analysis and Evaluation}}
	\small
	\textbf{Question:} What is the main purpose of a walkthrough?
	\begin{enumerate}[label=\Alph*.,itemsep=1pt,topsep=3pt]
			\item To confirm understanding of how a process and its controls work by following a transaction from initiation to recording.
			\item To conclude on operating effectiveness of all controls.
			\item To satisfy the process owner's request.
			\item To replace sampling.
	\end{enumerate}
\end{frame}
\begin{frame}[t]{Q87 -- Explanation}
	\footnotesize
	\textbf{Correct: A.} A walkthrough confirms design and understanding and informs later testing. \\[5pt]
\textbf{B wrong:} A single transaction cannot support such a conclusion. \\[3pt]
\textbf{C wrong:} Its purpose is audit understanding, not owner satisfaction. \\[3pt]
\textbf{D wrong:} Walkthroughs complement rather than substitute for testing. \\[3pt]
	\vfill {\tiny Part 2, Set 1}
\end{frame}
% ===== Q88 =====
\begin{frame}[t]{Q88 \hfill {\scriptsize Information Gathering, Analysis and Evaluation}}
	\small
	\textbf{Question:} ITGC testing shows user access management is ineffective. What is the effect on the engagement?
	\begin{enumerate}[label=\Alph*.,itemsep=1pt,topsep=3pt]
			\item Reliance on automated controls and system reports is reduced, so more direct testing of the data is needed.
			\item It can be dismissed as outside the scope.
			\item No effect if application controls passed a one-time test.
			\item It affects only the IT department.
	\end{enumerate}
\end{frame}
\begin{frame}[t]{Q88 -- Explanation}
	\footnotesize
	\textbf{Correct: A.} Weak access controls undermine the reliability of what the systems produce. \\[5pt]
\textbf{B wrong:} ITGC are a foundation for reliance on system-dependent evidence. \\[3pt]
\textbf{C wrong:} Controls could have been changed after the test. \\[3pt]
\textbf{D wrong:} System-dependent processes across the organization are affected. \\[3pt]
	\vfill {\tiny Part 2, Set 1}
\end{frame}
% ===== Q89 =====
\begin{frame}[t]{Q89 \hfill {\scriptsize Information Gathering, Analysis and Evaluation}}
	\small
	\textbf{Question:} To test revenue cutoff, which approach is BEST?
	\begin{enumerate}[label=\Alph*.,itemsep=1pt,topsep=3pt]
			\item Select invoices from the middle of the year.
			\item Compare revenue to budget.
			\item Select shipments dated shortly before and after period end and trace them to invoices and the period recorded.
			\item Review the revenue policy.
	\end{enumerate}
\end{frame}
\begin{frame}[t]{Q89 -- Explanation}
	\footnotesize
	\textbf{Correct: C.} Cutoff errors occur around period end, so testing focuses there. \\[5pt]
\textbf{A wrong:} Mid-year items do not test period-end timing. \\[3pt]
\textbf{B wrong:} Budget comparison is not a cutoff test. \\[3pt]
\textbf{D wrong:} The policy shows design, not timing in practice. \\[3pt]
	\vfill {\tiny Part 2, Set 1}
\end{frame}
% ===== Q90 =====
\begin{frame}[t]{Q90 \hfill {\scriptsize Information Gathering, Analysis and Evaluation}}
	\small
	\textbf{Question:} At an inventory count, the auditor selects items from the perpetual records and locates them on the floor. What does this test?
	\begin{enumerate}[label=\Alph*.,itemsep=1pt,topsep=3pt]
			\item Existence of recorded inventory.
			\item Completeness of recorded inventory.
			\item Valuation of inventory.
			\item Rights and obligations.
	\end{enumerate}
\end{frame}
\begin{frame}[t]{Q90 -- Explanation}
	\footnotesize
	\textbf{Correct: A.} Records to floor tests whether recorded items exist; floor to records tests completeness. \\[5pt]
\textbf{B wrong:} Completeness is tested by selecting items on the floor and tracing them to records. \\[3pt]
\textbf{C wrong:} Valuation concerns the amounts used, not presence. \\[3pt]
\textbf{D wrong:} Rights concern ownership, not physical presence. \\[3pt]
	\vfill {\tiny Part 2, Set 1}
\end{frame}

%%%%%%%%%% Engagement Supervision and Communication %%%%%%%%%%
% ===== Q91 =====
\begin{frame}[t]{Q91 \hfill {\scriptsize Engagement Supervision and Communication}}
	\small
	\textbf{Question:} A workpaper file shows reviewer sign-offs all dated on one day three weeks after fieldwork ended, with no review notes. What is the MOST supportable conclusion?
	\begin{enumerate}[label=\Alph*.,itemsep=1pt,topsep=3pt]
			\item The file does not show supervision during the engagement, which should be ongoing so issues are corrected in time.
			\item Late review is acceptable for low-risk engagements.
			\item Review notes are optional if sign-offs exist.
			\item Supervision is adequate because all sign-offs are present.
	\end{enumerate}
\end{frame}
\begin{frame}[t]{Q91 -- Explanation}
	\footnotesize
	\textbf{Correct: A.} Supervision should occur throughout the engagement and be evidenced. \\[5pt]
\textbf{B wrong:} Risk affects depth of review, not whether it happens during the work. \\[3pt]
\textbf{C wrong:} Evidence of review quality is expected. \\[3pt]
\textbf{D wrong:} Sign-offs without timing or notes do not show real review. \\[3pt]
	\vfill {\tiny Part 2, Set 1}
\end{frame}
% ===== Q92 =====
\begin{frame}[t]{Q92 \hfill {\scriptsize Engagement Supervision and Communication}}
	\small
	\textbf{Question:} A reviewer finds a conclusion in a draft workpaper that the evidence does not support. What is the MOST appropriate action?
	\begin{enumerate}[label=\Alph*.,itemsep=1pt,topsep=3pt]
			\item Approve it and note a follow-up for next year.
			\item Leave it as is since the auditor is experienced.
			\item Delete the conclusion silently.
			\item Return it for more evidence or a revised conclusion and document the resolution.
	\end{enumerate}
\end{frame}
\begin{frame}[t]{Q92 -- Explanation}
	\footnotesize
	\textbf{Correct: D.} Conclusions must be supported by evidence, and the review resolution must be evident. \\[5pt]
\textbf{A wrong:} Deferring the issue allows unsupported conclusions to stand. \\[3pt]
\textbf{B wrong:} Experience does not replace evidence. \\[3pt]
\textbf{C wrong:} Silent deletion hides a quality problem. \\[3pt]
	\vfill {\tiny Part 2, Set 1}
\end{frame}
% ===== Q93 =====
\begin{frame}[t]{Q93 \hfill {\scriptsize Engagement Supervision and Communication}}
	\small
	\textbf{Question:} A senior manager asks an auditor to remove an unfavorable finding from a workpaper. What should the auditor do?
	\begin{enumerate}[label=\Alph*.,itemsep=1pt,topsep=3pt]
			\item Delete the entire workpaper.
			\item Remove it and keep a private copy.
			\item Remove it if the manager provides a reason.
			\item Decline, document the request, and report it through the supervisor, CAE, or appropriate channel.
	\end{enumerate}
\end{frame}
\begin{frame}[t]{Q93 -- Explanation}
	\footnotesize
	\textbf{Correct: D.} Altering results at a manager's request compromises integrity and objectivity. \\[5pt]
\textbf{A wrong:} Deleting the workpaper destroys evidence. \\[3pt]
\textbf{B wrong:} A private copy does not cure an improper alteration. \\[3pt]
\textbf{C wrong:} A reason does not justify altering evidence-based results. \\[3pt]
	\vfill {\tiny Part 2, Set 1}
\end{frame}
% ===== Q94 =====
\begin{frame}[t]{Q94 \hfill {\scriptsize Engagement Supervision and Communication}}
	\small
	\textbf{Question:} An auditor is asked to sign off a workpaper prepared by a colleague without having reviewed it. What is the BEST response?
	\begin{enumerate}[label=\Alph*.,itemsep=1pt,topsep=3pt]
			\item Decline until a proper review is done, and document the reason.
			\item Sign to keep the engagement on schedule.
			\item Sign and add a note saying it was not reviewed.
			\item Alter the workpaper to match the auditor's own view.
	\end{enumerate}
\end{frame}
\begin{frame}[t]{Q94 -- Explanation}
	\footnotesize
	\textbf{Correct: A.} A sign-off represents that work has been reviewed. \\[5pt]
\textbf{B wrong:} Signing without review misrepresents the supervision performed. \\[3pt]
\textbf{C wrong:} A caveat does not make a sign-off truthful. \\[3pt]
\textbf{D wrong:} Altering another person's work is improper. \\[3pt]
	\vfill {\tiny Part 2, Set 1}
\end{frame}
% ===== Q95 =====
\begin{frame}[t]{Q95 \hfill {\scriptsize Engagement Supervision and Communication}}
	\small
	\textbf{Question:} After an engagement, how should the supervisor BEST evaluate an auditor's performance?
	\begin{enumerate}[label=\Alph*.,itemsep=1pt,topsep=3pt]
			\item Base it on hours charged.
			\item Base it only on the client's satisfaction survey.
			\item Wait for the annual appraisal.
			\item Give timely feedback against defined competencies and engagement quality criteria, and document it.
	\end{enumerate}
\end{frame}
\begin{frame}[t]{Q95 -- Explanation}
	\footnotesize
	\textbf{Correct: D.} Evaluation supports improvement and quality through timely, criteria-based feedback. \\[5pt]
\textbf{A wrong:} Hours do not measure quality. \\[3pt]
\textbf{B wrong:} Client satisfaction does not measure audit quality or independence. \\[3pt]
\textbf{C wrong:} Waiting delays learning and correction. \\[3pt]
	\vfill {\tiny Part 2, Set 1}
\end{frame}
% ===== Q96 =====
\begin{frame}[t]{Q96 \hfill {\scriptsize Engagement Supervision and Communication}}
	\small
	\textbf{Question:} During fieldwork the auditor finds a significant control deficiency. When should it be communicated?
	\begin{enumerate}[label=\Alph*.,itemsep=1pt,topsep=3pt]
			\item Informally to junior staff of the auditee.
			\item Promptly to the supervisor and appropriate management, rather than waiting for the final report.
			\item Only in the final report.
			\item At the closing meeting.
	\end{enumerate}
\end{frame}
\begin{frame}[t]{Q96 -- Explanation}
	\footnotesize
	\textbf{Correct: B.} Significant matters should be communicated in time for management to respond and for supervisors to adjust the work. \\[5pt]
\textbf{A wrong:} Junior staff may lack the authority to act. \\[3pt]
\textbf{C wrong:} The final report is too late for matters needing quick action. \\[3pt]
\textbf{D wrong:} Waiting until the close delays corrective action. \\[3pt]
	\vfill {\tiny Part 2, Set 1}
\end{frame}
% ===== Q97 =====
\begin{frame}[t]{Q97 \hfill {\scriptsize Engagement Supervision and Communication}}
	\small
	\textbf{Question:} A draft observation uses accusatory language and omits management's relevant explanation. Which communication quality criteria are MOST affected?
	\begin{enumerate}[label=\Alph*.,itemsep=1pt,topsep=3pt]
			\item Constructive and concise.
			\item Concise and timely.
			\item Clear and accurate.
			\item Objective and complete.
	\end{enumerate}
\end{frame}
\begin{frame}[t]{Q97 -- Explanation}
	\footnotesize
	\textbf{Correct: D.} Accusatory wording undermines objectivity and omitting relevant explanations affects completeness. \\[5pt]
\textbf{A wrong:} Concision is not the problem here. \\[3pt]
\textbf{B wrong:} Neither concision nor timeliness is the main issue. \\[3pt]
\textbf{C wrong:} Clarity and accuracy may be intact, while fairness is not. \\[3pt]
	\vfill {\tiny Part 2, Set 1}
\end{frame}
% ===== Q98 =====
\begin{frame}[t]{Q98 \hfill {\scriptsize Engagement Supervision and Communication}}
	\small
	\textbf{Question:} The engagement client has delayed providing data for three weeks. What is the BEST action?
	\begin{enumerate}[label=\Alph*.,itemsep=1pt,topsep=3pt]
			\item Conclude based on available data without stating the limitation.
			\item Wait indefinitely.
			\item Contact the board directly, bypassing the CAE.
			\item Escalate through the supervisor according to the escalation protocol, and document the effect on schedule and scope.
	\end{enumerate}
\end{frame}
\begin{frame}[t]{Q98 -- Explanation}
	\footnotesize
	\textbf{Correct: D.} Delays that affect the engagement need escalation and documentation. \\[5pt]
\textbf{A wrong:} A conclusion that hides a limitation is misleading. \\[3pt]
\textbf{B wrong:} Waiting allows the delay to threaten the engagement. \\[3pt]
\textbf{C wrong:} Escalation follows the CAE and agreed protocol. \\[3pt]
	\vfill {\tiny Part 2, Set 1}
\end{frame}
% ===== Q99 =====
\begin{frame}[t]{Q99 \hfill {\scriptsize Engagement Supervision and Communication}}
	\small
	\textbf{Question:} Fieldwork finds a critical, unpatched, internet-facing server. How should this be communicated?
	\begin{enumerate}[label=\Alph*.,itemsep=1pt,topsep=3pt]
			\item On a shared portal accessible to all auditees.
			\item In the weekly status e-mail to all stakeholders.
			\item In the final report only.
			\item Immediately by a secure channel to IT security leadership with limited distribution, and follow up in writing.
	\end{enumerate}
\end{frame}
\begin{frame}[t]{Q99 -- Explanation}
	\footnotesize
	\textbf{Correct: D.} Sensitive, urgent matters need prompt, controlled communication. \\[5pt]
\textbf{A wrong:} A widely accessible location exposes sensitive details. \\[3pt]
\textbf{B wrong:} Broad distribution of vulnerability details creates additional risk. \\[3pt]
\textbf{C wrong:} Waiting for the final report is too late. \\[3pt]
	\vfill {\tiny Part 2, Set 1}
\end{frame}
% ===== Q100 =====
\begin{frame}[t]{Q100 \hfill {\scriptsize Engagement Supervision and Communication}}
	\small
	\textbf{Question:} At the closing meeting, the process owner disputes a factual statement in a draft observation. What is the BEST approach?
	\begin{enumerate}[label=\Alph*.,itemsep=1pt,topsep=3pt]
			\item Escalate the dispute to the board immediately.
			\item Evaluate any new evidence, revise if warranted, and document the resolution; retain the observation if the evidence supports it.
			\item Refuse to consider any additional evidence.
			\item Delete the observation to preserve the relationship.
	\end{enumerate}
\end{frame}
\begin{frame}[t]{Q100 -- Explanation}
	\footnotesize
	\textbf{Correct: B.} Facts should be confirmed objectively, and observations stand or change on evidence. \\[5pt]
\textbf{A wrong:} Escalation is premature before the evidence is evaluated. \\[3pt]
\textbf{C wrong:} Ignoring evidence undermines accuracy and fairness. \\[3pt]
\textbf{D wrong:} Relationships should not drive changes to evidence-based findings. \\[3pt]
	\vfill {\tiny Part 2, Set 1}
\end{frame}
\begin{frame}{Answer key}
\small
1:C \quad 2:A \quad 3:D \quad 4:C \quad 5:C \quad 6:D \quad 7:B \quad 8:D \quad 9:B \quad 10:B \quad 11:C \quad 12:B \quad 13:A \quad 14:D \quad 15:A \quad 16:C \quad 17:D \quad 18:B \quad 19:C \quad 20:C \quad 21:B \quad 22:B \quad 23:D \quad 24:A \quad 25:A \\[6pt]
26:C \quad 27:C \quad 28:D \quad 29:C \quad 30:C \quad 31:D \quad 32:B \quad 33:B \quad 34:B \quad 35:C \quad 36:A \quad 37:D \quad 38:C \quad 39:A \quad 40:D \quad 41:A \quad 42:C \quad 43:A \quad 44:C \quad 45:C \quad 46:B \quad 47:A \quad 48:C \quad 49:D \quad 50:A \\[6pt]
51:B \quad 52:B \quad 53:D \quad 54:C \quad 55:B \quad 56:C \quad 57:D \quad 58:B \quad 59:C \quad 60:A \quad 61:D \quad 62:D \quad 63:A \quad 64:B \quad 65:B \quad 66:A \quad 67:B \quad 68:B \quad 69:D \quad 70:C \quad 71:B \quad 72:A \quad 73:D \quad 74:A \quad 75:D \\[6pt]
76:B \quad 77:A \quad 78:B \quad 79:C \quad 80:A \quad 81:B \quad 82:A \quad 83:D \quad 84:C \quad 85:C \quad 86:A \quad 87:A \quad 88:A \quad 89:C \quad 90:A \quad 91:A \quad 92:D \quad 93:D \quad 94:A \quad 95:D \quad 96:B \quad 97:D \quad 98:D \quad 99:D \quad 100:B
\end{frame}
\end{document}
